\chapter{GPU、显存与显示系统}

\begin{keyidea}
GPU 用大量并行执行资源换取吞吐，并以命令队列、显存和 fence 接入整机。本章从执行模型一直走到显示扫描边界。
\end{keyidea}

CPU 为单条指令的时延优化，GPU 为海量数据的吞吐优化——这是理解 GPU 的全部出发点。为了吞吐，GPU 把芯片面积让给成排的算术单元，把单线程性能让给了同时运行的上万个线程；为了喂饱这些单元，它配了独立的显存和巨大的带宽。但 GPU 对系统程序员来说并不神秘：命令照样走队列，数据照样走 DMA，完成照样走 fence 和中断。本章把 GPU 从“显卡”还原成一台吞吐优先的并行处理器和它的显示子系统。

\begin{longtable}{L{0.2\linewidth}L{0.38\linewidth}L{0.32\linewidth}}
\toprule
对象 & 与 CPU 的差别 & 本章回答的问题 \\
\midrule
执行模型 & 上万线程批量推进，不追单线程时延 & SIMT 怎样组织并行 \\\tablerowrule
显存 & 独立 VRAM/HBM，带宽数倍于主存 & 数据怎样进出 GPU \\\tablerowrule
命令路径 & 命令环+DMA+fence，不逐条同步 & CPU 与 GPU 怎样协作而不互等 \\\tablerowrule
显示子系统 & 按固定时序扫描帧缓冲 & 渲染完成为什么不等于显示完成 \\
\bottomrule
\end{longtable}

\begin{keyidea}
GPU 的性能秘密不是“核多”，而是“用海量在途工作隐藏一切时延”。读 GPU 文档时，先问在途有多少线程、数据在谁的存储里、同步在哪个边界。
\end{keyidea}

\section{一、GPU 为什么长得不像 CPU}

本节从一块硅片的预算说起。可用的晶体管总数大体
固定，怎么花，决定了芯片擅长什么：CPU 把预算的
大头给了控制与缓存，因为通用程序里单条指令等不
起；GPU 把预算给了成排的算术单元，因为它要处理
的工作天然是上百万份互不依赖的小任务。

\topic{时延优先与吞吐优先的面积分配}
芯片设计是一笔零和的账：功耗与面积的预算固定，
多给控制逻辑一平方毫米，算术单元就少一排。两种
处理器的分歧，从预算表的第一行就开始了。

CPU 的预算花在并不直接算数的地方。一个现代乱序
核里，真正执行运算的单元——整数 ALU、浮点单元
——只占核面积的一小部分，大约一两成；剩下的大
部分给了三件事。分支预测器赌下一条指令在哪里，
赌错了整条流水线清空重来；乱序调度维护几十到上
百条在途指令的重排序缓冲、重命名寄存器与发射队
列，只为让一条指令流不必等前面那条慢指令；大容
量 cache 则把平均访存时延从几百拍压到几拍。这
三件事有一个共同的客户：单线程。它们服务的不是
“多算”，而是“这条指令流别停”。

GPU 把同一笔预算反过来花。一个 SM（streaming
multiprocessor）里，取指、译
码、发射只做一套，却同时驱动 32 条算术道：控制
面积被 32 路摊薄，省下的预算再复制成几十上百个
这样的 SM。分支预测很简陋，乱序完全没有，cache
也不大——在 GPU 的任务假设下，这些都是浪费：
既然片上随时驻留着上万个互不依赖的线程，又何必
为某一个线程的去向破费晶体管？

两种机器的超标量也值得一比。CPU 的超标量是一条
指令流每拍发射多条——高端核做到 4 到 6 条——
代价是每拍都要为候选指令独立检查相关性、争用端
口；GPU 是一条指令带 32 份数据，相关性整组只检
查一次。同样追求“一拍多条”，一个把复杂度花在
“找”，一个花在“排”。设一套控制的面积与一条算
术道相当，那么“一套控制带一条道”的机器，控制
占去一半面积；“一套控制带 32 道”之后，控制占
比降到 $1/(1+32)\approx3\%$，同样的预算下算术
道接近翻倍。这就是 GPU 的面积账，第二节会看到
这笔账在执行时的形态。

把这笔预算账画成两块硅片的草图，两种押注一目了然：
同一方硅片上预算是零和的，控制与缓存多占一分，算
术道就少一排。

\begin{pdfdiagram}
  % 左：CPU —— 控制与缓存占大头
  \draw[BookRule, line width=0.5pt, dashed, rounded corners=4pt] (0,0) rectangle (4.7,4.6);
  \node[hw label] at (2.35,4.88) {CPU：时延优先};
  \node[hw control, minimum width=3.7cm, minimum height=1.0cm] at (2.35,3.85) {控制：预测·乱序·调度\\约三成};
  \node[hw memory, minimum width=3.7cm, minimum height=1.6cm] at (2.35,2.2) {大容量 cache\\约一半};
  \node[hw compute, minimum width=3.7cm, minimum height=0.7cm] at (2.35,0.65) {算术单元：约一两成};
  % 右：GPU —— 算术道占大头
  \draw[BookRule, line width=0.5pt, dashed, rounded corners=4pt] (5.7,0) rectangle (10.9,4.6);
  \node[hw label] at (8.3,4.88) {GPU：吞吐优先};
  \node[hw control, minimum width=2.4cm, minimum height=0.62cm] at (8.3,4.06) {一套控制带 32 道：约 3\%};
  \node[hw label] at (8.3,3.52) {成排算术道：面积的绝大部分};
  \foreach \r in {0,...,3} {
    \foreach \c in {0,...,5} {
      \node[hw compute, minimum width=0.62cm, minimum height=0.48cm, inner xsep=1pt, inner ysep=1pt] at ({6.35+0.8*\c},{3.05-0.62*\r}) {};
    }
  }
  \node[hw memory, minimum width=1.7cm, minimum height=0.5cm] at (6.95,0.48) {小 cache};
  \node[hw label] at (9.25,0.48) {时延靠线程隐藏};
\end{pdfdiagram}
\diagramnote{同一笔面积预算的两种花法（比例为示意）。CPU 的大头给了控制与缓存：它们不直接算数，只为单条指令流不停顿；GPU 一套发射逻辑带 32 条算术道，控制占比摊薄到约 $1/(1+32)\approx3\%$，省下的预算再复制成更多算术道。两种分法没有高下：一个押“这条指令流别停”，一个押“每秒总量最大”。}

顺着这笔账问到底：既然要并行，为什么不把同样的
预算堆成一百个乱序核？答案是汇率不划算。一个乱
序核连同它的预测器、调度窗口与私有 cache，面积
抵得上几十条算术道；同样一块硅片，堆乱序核得到
一两百个在途线程，堆 SM 得到二十六万。并行度的
汇率差了三个数量级，而 GPU 的工作假设——海量、
同构、互不依赖——恰恰用不上乱序核买来的那些单
线程能力。预算的每种花法，都是在赌自己的工作长
成什么样；GPU 赌赢了图形，图形也塑造了 GPU。

\begin{longtable}{L{0.2\linewidth}L{0.35\linewidth}L{0.35\linewidth}}
\toprule
预算去向 & CPU（时延优先） & GPU（吞吐优先） \\
\midrule
算术单元 & 核面积的小头，约一两成 & 面积的绝大部分，成排复制 \\\tablerowrule
控制逻辑 & 预测、乱序、重命名，占大头 & 一套发射带 32 道，能省则省 \\\tablerowrule
缓存 & MB 到上百 MB，压低平均时延 & 每 SM 小容量，时延靠线程藏 \\\tablerowrule
优化对象 & 单个任务的完成时间 & 每秒完成的总工作量 \\
\bottomrule
\end{longtable}

要防止一个误读：这不是“先进与落后”的对照。GPU
核的简陋是刻意为之——在“工作可拆成海量小任务”
的假设下，复杂的控制逻辑只增加功耗，不增加吞吐；
反过来，把 GPU 的小核搬去跑通用程序，会因为单
线程太慢而根本没法用。“冒险、异常与性能”一章把时延与吞吐拆成两
个独立的量，这里看到的是这两个量在硅片上的形状：
两种芯片各自把预算押给了其中一个。

这个分工也不是从来如此。1990 年代的“显卡”只是
帧缓冲加固定功能管线，2000 年代可编程着色器让算
术单元成片出现，再往后统一着色器架构把顶点与像
素单元合并成通用的算术道，SIMT 的形态就此定型。
芯片演化的每一步，都是把预算从固定功能挪给可编
程算术——本节的面积账，在三十年里被反复执行。

\begin{classicnote}
P\&H 讲 GPU 的那一节开门见山：GPU 是为吞吐优化
的多处理器，它要回答的不是“一条指令多快”，而是
“怎样让每秒数百亿次运算流过芯片”。本节后面每一
项结构差别，都是这句话的推论。
\end{classicnote}

\topic{图形需求天然是海量并行}
这种结构配得上什么工作？先看图形给出的数字。一
帧 1080p 画面有 $1920\times1080=2,073,600$ 个像
素，4K 画面有 $3840\times2160=8,294,400$ 个；按
60 帧每秒，GPU 每秒要着色的像素是 1.2 亿到 5 亿
个。每个像素的着色——纹理采样、光照、混合——
要几十到几百次运算，乘起来就是每秒 $10^{10}$ 到
$10^{12}$ 次运算的量级。顶点一侧同理：一个精细
角色就有几万个顶点，一屏上百个角色加场景，每帧
顶点数以千万计，每个顶点做一次 $4\times4$ 矩阵
变换约 32 次浮点运算，仅几何变换又是每秒
$10^{10}$ 次上下的一笔。

比数量更重要的是形状：这些计算几乎互不依赖。第
五十万个像素的颜色，不等第四十九万个算完；每个
顶点的变换，不需要知道相邻顶点的结果。这种“海
量、同构、互不依赖”的任务叫数据并行，口语里叫
“令人愉快的并行”——愉快之处就在于并行度是白
送的，不需要任何拆分技巧。实际负载只会更重：半
透明叠加与过绘制让同一像素被重复着色，每帧的片
元数往往是像素数的 2 到 4 倍，上面还是保守估计。

拿帧预算反推一次结构。60 帧每秒给每帧 16.7 ms。
假设每像素只花 100 拍，把 1080p 一帧全部串行放
在一条 2 GHz 的算术道上跑：

\begin{lstlisting}
串行时间 = 2,073,600 像素 x 100 拍 / 2 GHz
        ≈ 104 ms，是帧预算的 6 倍多
再乘两条：每像素实际远不止 100 拍，
        几何与光栅化还没算进去
结论：需要 10^2 到 10^3 条并行的算术道
\end{lstlisting}

帧预算还有第二层含义：平均值之外是稳定。固定节
奏的影片只要帧率恒定就流畅，交互图形的 60 帧却
怕抖动——某一帧多花了 5 ms，人立刻感到卡顿。
这对结构的要求是：并行度不仅要够，还要为最坏的
场景留余量，工程上常按峰值需求的一两倍配置算力。
批量小核在这里又赢一次：加并行度远比提单线程速
度便宜，余量因此给得起。

任务结构决定机器结构。工作是一百万个同构的小题
时，配一百万道简装工序，远比配一个全能大师傅划
算——大师傅再快，一次也只能做一道题。这就是
“批量小核”的必然性：不是工程师偏爱核多，是任务
特征决定了这种结构。章首核心思想说 GPU 的
秘密是“用海量在途工作隐藏一切时延”，那句话的
前提就在这里：需求本身得先提供足够多的在途工作。

这段推理还给出一张可复用的判定单：一个工作适配
GPU 的三个条件——任务量以百万计、每份同构、份
与份互不依赖。缺第一条，下发的开销收不回；缺第
二条，分支分歧降低吞吐；缺第三条，同步开销会主导总体性能。
存储与设备事务相关章节讲 DMA 时说过它只擅长“整块、连续、大批”
的搬运，这里看到的是同一哲学在计算上的版本：批
量换效率。

\topic{GPU 的基本组成}
一个 SM（AMD 阵营叫计算单元 CU）的内部，是把第
四章最小 CPU 的思路改
了比例：部件还是那些部件，配比完全不同。最小
CPU 里一套控制带一条算术道；SM 里一套发射带 32
条道，这样的组再复制四份。

\begin{longtable}{L{0.2\linewidth}L{0.28\linewidth}L{0.4\linewidth}}
\toprule
部件 & 规模量级（每 SM） & 职责 \\
\midrule
算术道 & 128 道 FP32，分 4 组 & 执行数值运算，每组每拍消化一条 warp 指令 \\\tablerowrule
寄存器堆 & 65,536 个 32 位，共 256 KiB & 全体驻留线程的上下文常驻其中 \\\tablerowrule
共享内存/L1 & 几十到一百多 KiB & 块内线程交换数据的近端存储 \\\tablerowrule
warp 调度器 & 4 个 & 每拍各挑一组就绪 warp 发一条指令 \\\tablerowrule
特殊功能单元 & 少量 & 超越函数、纹理采样等杂活 \\
\bottomrule
\end{longtable}

寄存器堆值得单独讲，因为它是“换线程零代价”的
物理基础。CPU 切换线程要把当前寄存器逐个保存到
内存、再装入下一个线程的，一次切换几百拍起步，
操作系统线程更是微秒级。GPU 的做法是让所有驻留
线程的上下文永远留在寄存器堆里：每个 warp 的 32
个线程各占一段，“切换”只是下一拍换一组来执行，
没有任何保存与恢复。代价就是那
$65,536\times4\ \text{B}=256$ KiB——比 SM 自己
的 L1 还大。GPU 实际上把别的芯片给 cache 的预算，
挪给了“驻留更多线程”；这与它的时延哲学一致：
不去缩短一次访存的时延，而是让这段时延里永远有
别的线程可算。调度器与寄存器堆怎么配合，第二节
讲占用率时还要回来算细账。

和处理器核心相关章节的寄存器组对照最直观。最小 CPU 里 8 到
16 个寄存器就够用，因为一次只执行一条指令，状
态随指令流动；SM 的 65,536 个寄存器按 32 线程一
组切成 warp 上下文，全部静态摊在明面上。GPU 片
上最大的单片存储不是 cache，是寄存器堆——它把
“谁在现场”做成静态分配，而不是动态换人。

特殊功能单元与纹理单元是图形时代的遗产：正弦、
开方、倒数这类超越函数用专用硬件算，比软件查表
快得多；纹理单元自带插值与寻址模式。通用计算兴
起后它们留了下来，只是换了服务的负载。这一节不
展开，记住 SM 里不只有算术道即可。

总算力可以把数字直接乘出来。一张现代中高端卡有
上百个 SM，每个 SM 每拍完成 128 个 FP32 乘加：

\begin{lstlisting}
总算力 = SM 数 x 每拍乘加数 x 2（乘加算两次运算）x 时钟
      = 128 x 128 x 2 x 1.5 GHz
      ≈ 49 TFLOPS
同代通用 CPU 约低一到两个数量级
\end{lstlisting}

这个数字单独看没有用处，配上带宽才完整——第十二
章 roofline 的拐点正是“峰值算力除以带宽”，带宽
一侧第三节补齐，这里先记住算力的量级。

\topic{GPU 与 CPU 的规格对照}
把两种芯片并排摆开，每一项差别都能回溯到前面的
面积分配：

\begin{longtable}{L{0.18\linewidth}L{0.36\linewidth}L{0.36\linewidth}}
\toprule
项目 & 桌面/服务器 CPU & 现代独立显卡 GPU \\
\midrule
在途线程数 & 核数乘超线程，16--128 & 上万，满配约 26 万 \\\tablerowrule
时钟 & 4--6 GHz & 1.5--2.5 GHz \\\tablerowrule
主存储带宽 & DDR5 双通道约 77 GB/s & GDDR6 数百 GB/s，HBM 过 1 TB/s \\\tablerowrule
cache 策略 & 大容量，缩短时延本身 & 小容量，靠线程数填满时延 \\\tablerowrule
单线程速度 & 基准 & 慢数倍 \\\tablerowrule
优化目标 & 单任务时延 & 全片总吞吐 \\
\bottomrule
\end{longtable}

逐项看。线程数：GPU 的在途数是 warp 数乘 32，按
$128\times64\times32=262,144$ 计，满配二十六万
线程同时在片上；CPU 满打满算一百出头。注意这是
“驻留”而不是“创建”：二十六万线程的上下文全
在寄存器堆里，随时可以被调度器点名。

时钟：GPU 的算术阵列太宽，频率抬不上去。“电源、时钟与信号完整性”一章
讲电源时算过，动态功耗随频率与电压一起涨，频率
抬一档、电压往往还要跟一档；把这笔账乘以上万条
算术道，就知道 GPU 的时钟为什么停在 CPU 的一半
上下。单条道因此更慢，但道数把差距赢回来几十倍。

带宽：DDR5 双通道按
$4800\ \text{MT/s}\times8\ \text{B}\times2\approx77$
GB/s；GDDR6 用 384 位宽接口，按
$16\ \text{Gbps}\times384/8=768$ GB/s；HBM 把存
储堆叠到处理器旁边，整卡过 1 TB/s。十倍的带宽差
距不是锦上添花，是 GPU 结构的另一半——第三节专
门讲这笔带宽从哪来、怎么花。

cache 策略的分歧最值得回味。CPU 用大 cache 把时
延本身缩短：平均一次访存从几百拍变成几拍。GPU
的 cache 小，因为它不打算缩短时延，而打算用时间
换：访存的几百拍里，调度器换别的 warp 上道，时
延被“填满”而不是被“缩短”。“冒险、异常与性能”一章说“用并行度
买吞吐、买不来时延”，两种芯片把这句话各自实现
了一遍，只是买的方向相反。

这张表的读法也因此明确：不要跨栏比较单核性能。
两栏优化的本来就不是同一个量，跨栏比就像问“集
装箱船和快艇哪个好”——答案取决于你要运什么。

\topic{为什么 GPU 不追求单线程快}
把上面几行规格合起来，单线程慢是结构性的：时钟
低一半，没有乱序和多发射，分支预测简陋，发射宽
度还要被一个 warp 的 32 个线程共享。粗算差距：
时钟差 2 到 3 倍；指令级并行再差数倍——乱序核
一拍能退出多条，SIMT 单线程只能等整组共享的那
一条发射。乘起来，同一段串行代码在 GPU 上慢 5
到 10 倍是正常区间。这个数字不必精确，记住方向
就够：差距是结构性的，不随驱动或编译器消失。

这也不是缺陷，更不需要修复。GPU 的优化目标从来
不是“这一个线程多久算完”，而是“每秒完成多少
工作”。那能不能给 GPU 也装上乱序和大 cache？
预算零和：装了就得砍算术道。乱序窗口换来的是单
线程时延，关闭的通道降低的是总吞吐——对一块为吞
吐而生的芯片，这是亏本交换。工程上没有“什么都
要”的芯片，只有“押哪个量”的芯片。

也要补一句平衡：两种取向并非泾渭分明。近年的
GPU，cache 一代比一代大，线程调度一代比一代灵
活；CPU 的向量单元也一代比一代宽。边界在移动，
优先级没有变：GPU 添的每一项“聪明”都要先回答
“它给吞吐带来多少”，CPU 添的每一点并行都要先
回答“它拖不拖单线程”。看懂预算表，就看懂了这
种互相借鉴却不互相变成的演化。

“冒险、异常与性能”一章的拆分在这里收尾：时延与吞吐是两个独立的
量，改善一个不必改善另一个。GPU 把单线程时延整
块让渡出去，换回二十六万在途线程的总吞吐。由此
得到一个直接的工程推论：串行段、低并行度的代码
应该留在 CPU 上，只有能拆出几万个小任务的部分
才值得下发；“冒险、异常与性能”一章的 Amdahl 定律同时提醒，串行
段的长短决定整个混合系统的加速上限。CPU 与 GPU
的分工由此而来，“命令路径：命令环、fence 与显示扫描”一节的命令路径，就是为这
种分工设计的协作机制。

\begin{warningbox}
“GPU 核多，所以什么都能加速”是最常见的误读。拆
不出几万个独立任务的工作上了 GPU 反而更慢：单线
程本来就慢数倍，下发与回收本身还有一笔开销（第
四节）。GPU 加速的前提是任务的形状，不是芯片的
名气。
\end{warningbox}

\section{二、SIMD/SIMT 执行模型}

第一节回答了“GPU 为什么长这样”，本节回答“这些
算术道怎样被指令流驱动”。答案分两层：编程视角
是一组标量线程各算各的，硬件视角是 32 条道共享
一条指令流。

\topic{SIMD 与 SIMT 的区别}
SIMD（single instruction, multiple data）是一条
指令同时操作多个数据：一条 AVX-512 加法取出两个
512 位向量寄存器，一拍完成 16 对 float 相加。并
行度写在指令集里——向量的宽度是指令的一部分，
从 SSE 的 128 位一路加到 512 位；程序员看到的是
“一个线程，手里握着很宽的寄存器”，分支与循环仍
是标量语义。SIMT（single instruction, multiple
threads）换了组织方式：程序写的是一组标量线程，
每个线程有自己的寄存器、自己的数据、逻辑上自己
的程序计数器；硬件把 32 个步调一致的线程绑成一
组，共用一次取指与发射。并行度不写在指令里，写
在程序结构里——线程号决定数据号。

\begin{longtable}{L{0.18\linewidth}L{0.36\linewidth}L{0.36\linewidth}}
\toprule
维度 & SIMD & SIMT \\
\midrule
并行写在哪 & 指令里，宽度属于指令集 & 程序结构里，线程数属于启动配置 \\\tablerowrule
寄存器 & 一组宽向量寄存器 & 每线程一套标量寄存器 \\\tablerowrule
程序计数器 & 一个 & 逻辑上每线程一个，硬件每组一个 \\\tablerowrule
分支 & 向量内无分支，分支走标量路 & 组内分歧时各路轮流走（本节后文） \\\tablerowrule
典型形态 & SSE、AVX-512、NEON & CUDA/OpenCL 内核 \\
\bottomrule
\end{longtable}

两种模型能算同一件事，差别在“谁来暴露并行”。
SIMD 要编译器或程序员把数据亲手装进向量，分支一
多就装不满；SIMT 让程序员只管写一个线程的逻辑，
宽度的利用交给硬件调度。后者在分支多、访存模式
不规则的真实图形与计算代码里好写得多——这是
GPU 选 SIMT 的工程理由，不是理论上的必需。

GPU 的“线程”必须和操作系统线程分开称量。操作
系统线程创建要分配栈与内核结构，10 到 100 微秒
起步，切换一次也要微秒级；GPU 线程的“创建”只
是调度器在寄存器堆里划出一段、指派一条算术道，
切换是常驻换组，一拍完成，约合半纳秒。两端相差
三到四个数量级，所以 GPU 程序一次启动敢铺几十
万线程，服务器上的线程池却以几十计。第一节说寄
存器堆是“换线程零代价”的物理基础，这里看到的
是它的程序员视角：线程创建成本足够低，可按数据规模配置，
开线程这件事本身从预算表里消失了。

\begin{classicnote}
P\&H 把 SIMT 称为硬件视角的 SIMD：对程序员它是
一簇独立线程，对数据通路它就是向量。更早的血统
是向量机——一条指令定义为对一整串数据的循环；
GPU 相当于把这个循环摊开成 32 条并排的道。
\end{classicnote}

\topic{warp/wavefront 的调度}
SIMT 的执行单位叫 warp（NVIDIA，32 线程）或
wavefront（AMD，64 线程，新架构也可按 32 配置）。
一个 warp 一个程序计数器：调度器选中它，就发射
一条指令，组内 32 个线程的算术道同时执行各自的
数据。取指、译码、发射只做一次，服务 32 个线程
——第一节说的“控制面积摊薄”，指的就是这个结
构。

关键规则是：warp 是调度的最小单位，组内共进退。
只要有一个线程的操作数没准备好——典型情况是访
存的数据还没回来——整组这条指令就不能发射，调
度器把发射槽让给另一组就绪的 warp。实现上每个
warp 带着就绪标志，操作数齐了才置位；乱序核用
重排序缓冲在一团指令里找就绪的那条，SM 用整组
的就绪位在几十组 warp 里找就绪的那组——同一个
问题，两种规模的答案。

让出的那几百拍是这笔买卖的核心。一次显存访问要
四到八百拍（1.5 GHz 下约合 270 到 530 纳秒，
正是 GDDR 与 HBM 的真实量级），这段空窗不由发起
访存的 warp 硬等，而由别的 warp 的计算填满。这
就是 GPU 时延隐藏的第一机制：不缩短时延，用别
人的工作盖住时延。一个 SM 四个调度器每拍最多发
四条 warp 指令，即 128 个线程运算——第一节那
49 TFLOPS 就是这么一拍一拍垫出来的。

把这次让出与顶上画在时间轴上，“不缩短时延、用别
人的工作盖住时延”就有了具体的形状。

\begin{waveform}[x=0.8cm,y=0.85cm]
  % 参考线：warp0 发起访存 / 数据返回
  \draw[BookRule, line width=0.45pt, dashed] (3,4.15) -- (3,0.3);
  \draw[BookRule, line width=0.45pt, dashed] (10,4.15) -- (10,2.75);
  % warp 0
  \node[hw label, anchor=east] at (0.55,3.3) {warp 0};
  \draw[BookTeal, fill=BookCodeBg, line width=0.8pt] (1,3.0) rectangle (3,3.6);
  \node at (2,3.3) {执行};
  \draw[BookRule, dashed, line width=0.7pt] (3,3.0) rectangle (10,3.6);
  \node[hw label] at (6.5,3.3) {访存等待（让出发射槽）};
  \draw[BookTeal, fill=BookCodeBg, line width=0.8pt] (10,3.0) rectangle (12,3.6);
  \node at (11,3.3) {继续};
  % warp 1
  \node[hw label, anchor=east] at (0.55,2.1) {warp 1};
  \draw[BookRule, dashed, line width=0.7pt] (1,1.8) rectangle (3,2.4);
  \node[hw label] at (2,2.1) {驻留};
  \draw[BookTeal, fill=BookCodeBg, line width=0.8pt] (3,1.8) rectangle (6,2.4);
  \node at (4.5,2.1) {执行（顶上）};
  \draw[BookRule, dashed, line width=0.7pt] (6,1.8) rectangle (12,2.4);
  \node[hw label] at (9,2.1) {访存等待};
  % 算术道占用
  \node[hw label, anchor=east] at (0.55,0.9) {算术道};
  \draw[BookAccent, fill=BookCodeBg, line width=0.8pt] (1,0.6) rectangle (3,1.2);
  \node at (2,0.9) {warp 0};
  \draw[BookAccent, fill=BookCodeBg, line width=0.8pt] (3,0.6) rectangle (6,1.2);
  \node at (4.5,0.9) {warp 1};
  \draw[BookAccent, fill=BookCodeBg, line width=0.8pt] (6,0.6) rectangle (9,1.2);
  \node at (7.5,0.9) {warp 2};
  \draw[BookAccent, fill=BookCodeBg, line width=0.8pt] (9,0.6) rectangle (12,1.2);
  \node at (10.5,0.9) {warp 3};
  % 标注
  \draw[BookMuted, line width=0.6pt, <->] (3,4.42) -- node[hw label, above]{访存时延 400--800 拍} (10,4.42);
  \node[hw label] at (6.5,0.05) {时延没有被缩短，而是被在途 warp 的计算填满};
\end{waveform}
\diagramnote{warp 调度怎样隐藏访存时延：warp 0 发起访存后整组让出发射槽，调度器立刻换上驻留待派的 warp 1（以及更多 warp），算术道一行始终有主；数百拍后数据返回，warp 0 重新就绪继续。空窗由别人的工作盖住——这就是“不缩短时延、填满时延”的时序形状。}

这个思想本书已经见过一次。存储与外设队列案例讲 NVMe 队列
时算过：设备时延固定，吞吐靠在途请求数撑起来，
是“冒险、异常与性能”一章 Little 定律的直接应用。GPU 把同一公式
搬进处理器内部：访存时延固定，算术单元的利用率
靠在途 warp 数撑起来。一个叫队列深度，一个叫驻
留线程数，名字不同，机制是同一个。

顺带回答一个常见疑问：为什么不把 warp 定成 128
线程，把控制摊得更薄？因为分歧与合并访存的代价
都随组变大：组越大，越难凑齐同方向的分支，越难
凑出连续的一片地址，一次让出浪费的算术道也越多。
32 是“摊薄控制”与“组内一致性”之间的平衡点；
AMD 历史上用 64，新架构也退回可按 32 配置，原因
相同。

\topic{占用率：在途线程数决定隐藏能力}
调度器手里有多少 warp 可换，直接决定能盖住多长
的时延。这个量的行话叫占用率（occupancy）：

\begin{lstlisting}
占用率 = 本 SM 实际驻留的 warp 数 / 硬件上限
硬件上限的典型值是 64 warp/SM
\end{lstlisting}

需要多少才够，可以算。设一次访存时延 400 拍，
每个线程平均算 100 拍就要发起下一次访存：

\begin{lstlisting}
一个 warp 每 100 拍供给 100 拍工作，然后空等 400 拍
要算术道不停：至少需要 400/100 = 4 组 warp 轮转
计入发射间隔与毛刺，工程上取 8 组以上才稳
占用率低于 4/64 ≈ 6% 时，算术道大量空转
\end{lstlisting}

驻留上限本身由三样资源卡：寄存器堆、共享内存、
warp 槽位，三者取最紧的那个。寄存器堆的账最直
观：

\begin{lstlisting}
每线程 40 个寄存器：一个 warp 占 32 x 40 = 1,280 个
  65,536 / 1,280 = 51，可驻留 51 组 warp
每线程 255 个寄存器：一个 warp 占 8,160 个
  65,536 / 8,160 ≈ 8，只能驻留 8 组
  占用率 8/64 = 12.5%，400 拍的访存藏不干净
\end{lstlisting}

把这条计算链画成流程，编译器报告的寄存器数怎样一路决定占用率，
就一目了然：

\begin{pdfdiagram}
  % 例一：每线程 40 个寄存器
  \node[hw state, minimum width=2.7cm] (a1) at (1.6,0) {每线程 40 个寄存器\\每 warp 占 $32\times40=1{,}280$ 个};
  \node[hw compute, minimum width=2.5cm] (a2) at (4.75,0) {$65{,}536\div1{,}280=51$\\可驻留 51 组 warp};
  \node[hw compute, minimum width=2.2cm] (a3) at (7.5,0) {占用率\\$51/64\approx80\%$};
  \node[hw state, minimum width=1.9cm] (a4) at (10.0,0) {藏得住\\400 拍};
  \draw[hw bus] (a1.east) -- (a2.west);
  \draw[hw bus] (a2.east) -- (a3.west);
  \draw[hw bus] (a3.east) -- (a4.west);
  % 例二：每线程 255 个寄存器
  \node[hw state, minimum width=2.7cm] (b1) at (1.6,-1.7) {每线程 255 个寄存器\\每 warp 占 $8{,}160$ 个};
  \node[hw compute, minimum width=2.5cm] (b2) at (4.75,-1.7) {$65{,}536\div8{,}160\approx8$\\只能驻留 8 组};
  \node[hw compute, minimum width=2.2cm] (b3) at (7.5,-1.7) {占用率\\$8/64=12.5\%$};
  \node[hw control, minimum width=1.9cm] (b4) at (10.0,-1.7) {400 拍\\藏不住};
  \draw[hw bus] (b1.east) -- (b2.west);
  \draw[hw bus] (b2.east) -- (b3.west);
  \draw[hw bus] (b3.east) -- (b4.west);
  \node[hw label] at (5.5,-2.85) {分母固定：每 SM 65,536 个寄存器（256 KiB）；内核少用寄存器，是在给时延隐藏能力腾地方};
\end{pdfdiagram}
\diagramnote{寄存器预算到占用率的计算链（分母固定为每 SM 65,536 个 32 位寄存器，共 256 KiB）。每线程 40 个寄存器时一个 warp 占 1,280 个，可驻留 51 组，占用率约 80\%，400 拍的访存时延藏得住；每线程 255 个时一个 warp 占 8,160 个，只能驻留 8 组，占用率 12.5\%，发射槽开始空转。编译器报告每线程寄存器数，报的其实是占用率的上限。}

工程含义有两层。其一，内核少用寄存器不是洁癖，
是在给时延隐藏能力腾地方：编译器报告每线程寄存
器数时，报的其实是占用率的上限；为省寄存器把变
量溢出到内存也是常见交换——多花一次访存，换更
多在途 warp，只要换来的 warp 数盖得住新的访存就
划算。其二，占用率不是唯一杠杆：每线程若能多排
几条互相独立的指令（指令级并行），等待期间本组
warp 自己也有活干，4 组的下限就可以再低。第一
节说“时延被填满而不是被缩短”，填满它的可以是
更多的线程，也可以是每线程更多的独立指令——两
者相乘，才是总隐藏能力。

共享内存卡上限时，账法相同。设 SM 给共享内存的
配额是 100 KB，每个块申请 32 KB、含 8 组 warp：

\begin{lstlisting}
100 / 32 = 3，只能驻留 3 块，共 24 组 warp
占用率 24/64 = 37.5%
每块申请砍到 16 KB：驻留 6 块，共 48 组 warp
\end{lstlisting}

共享内存是块内线程的协作空间，用量越大单块越强，
但驻留越少——又一次“每线程便利”与“在途数量”
的交换。它具体怎么用，第三节讲显存层次时展开。

还要看清占用率与性能的关系不是线性的。从 4 组加
到 16 组，算术道的利用率显著提升，因为 400 拍的
空窗还没填满；从 48 组加到 64 组，空窗早已填满，
多出来的 warp 只是排队。隐藏时延这件事有饱和点，
调占用率的目标是“过拐点”，不是“拉满”——超过
拐点之后，该看的就不是占用率，而是访存模式与分
支，也就是接下来两笔账。

\topic{分支分歧的代价}
warp 共享一个程序计数器，遇到分支怎么办？硬件
的答案是：都走，轮流走。if-else 两边各有线程要
走时，先执行 A 侧、把走 B 的线程屏蔽成不活跃，
再执行 B 侧、屏蔽走 A 的，走完两边在汇合点重新
对齐。被屏蔽线程的算术道空转，这段空转就是分歧
的代价，行话叫分支分歧（branch divergence）。

代价可以精确算。设 A、B 两侧各花 $t$ 拍：

\begin{lstlisting}
无分歧：32 线程同走一路，warp 花 t 拍，吞吐 32/t
一半走 A 一半走 B：warp 花 2t 拍，
  期间每侧只有 16 线程活跃
  有效吞吐 = 32/(2t) = 16/t，直接减半
\end{lstlisting}

把这段过程画在时间轴上，哪一半线程在执行、哪一
半在空等，一眼可见。

\begin{waveform}[x=0.8cm,y=0.85cm]
  % 参考线：分支点 / A 结束 B 开始 / 汇合
  \draw[BookRule, line width=0.45pt, dashed] (1,3.85) -- (1,0.9);
  \draw[BookRule, line width=0.45pt, dashed] (5.5,3.85) -- (5.5,0.9);
  \draw[BookRule, line width=0.45pt, dashed] (10,3.85) -- (10,0.9);
  % 线程 0--15：先走 A，后被屏蔽
  \node[hw label, anchor=east] at (0.55,2.9) {线程 0--15};
  \draw[BookTeal, fill=BookCodeBg, line width=0.8pt] (1,2.6) rectangle (5.5,3.2);
  \node at (3.25,2.9) {路径 A 执行};
  \draw[BookRule, dashed, line width=0.7pt] (5.5,2.6) rectangle (10,3.2);
  \node[hw label] at (7.75,2.9) {屏蔽（算术道空转）};
  \draw[BookTeal, fill=BookCodeBg, line width=0.8pt] (10,2.6) rectangle (11.6,3.2);
  % 线程 16--31：先被屏蔽，后走 B
  \node[hw label, anchor=east] at (0.55,1.6) {线程 16--31};
  \draw[BookRule, dashed, line width=0.7pt] (1,1.3) rectangle (5.5,1.9);
  \node[hw label] at (3.25,1.6) {屏蔽（算术道空转）};
  \draw[BookTeal, fill=BookCodeBg, line width=0.8pt] (5.5,1.3) rectangle (10,1.9);
  \node at (7.75,1.6) {路径 B 执行};
  \draw[BookTeal, fill=BookCodeBg, line width=0.8pt] (10,1.3) rectangle (11.6,1.9);
  \node[hw label, anchor=west] at (10.1,3.5) {汇合：重新对齐};
  % 标注
  \draw[BookMuted, line width=0.6pt, <->] (1,4.12) -- node[hw label, above]{分歧区：A、B 串行，共 $2t$} (10,4.12);
  \node[hw label] at (5.5,0.55) {任一时刻只有一半算术道活跃 $\Rightarrow$ 有效吞吐减半};
\end{waveform}
\diagramnote{分支分歧的时间账：warp 共享一个程序计数器，A、B 两侧只能轮流执行——走 A 时 B 侧 16 条线程被屏蔽，走 B 时 A 侧被屏蔽，汇合点之后才恢复整组同行。两段路径的耗时相加，期间算术道始终有一半在空转，有效吞吐因此减半。}

嵌套起来更狠：每多一层分歧，顺序执行的路径数最
多翻一倍，$d$ 层嵌套最坏 $2^{d}$ 条路径挨个走；
32 个线程完全散开时，warp 退化成 32 条单线程路
径，吞吐只剩 1/32。较新的架构引入独立线程调度，
允许组内更灵活地分裂与汇合，但“同组同方向最便
宜”这条代价曲线没有变，只是毛刺少了。

工程含义与存储与设备事务相关章节 cache 友好遍历是同一类动作：按
机器的偏好排数据。分支前把同方向的数据聚到同一
个 warp（排序或分块），或者干脆把分支改写成无
分支的掩码算术——两侧都算、按条件择一，编译器
对短分支常自动做这个改写，长分支只能靠人。判断
该不该改写的尺子就是上面那笔账：两侧都短时，掩
码算术花 $2t$ 买无分歧；一侧明显更长时，先重排
数据再分歧。

\topic{合并访存}
warp 的一条 load 指令携带 32 个地址。内存子系统
不逐地址处理，而是先看这 32 个地址落在几条 128
字节段（cache 行）里，每条段发一笔事务。最好的
情况是连续线程访问连续地址：32 个线程各取 4 字
节，$32\times4=128$ 字节恰好填满一条段，拼成一
笔大事务，一次拿满，这叫合并访存（coalesced
access）。最坏的情况是跨步：每线程隔 32 个
float（128 字节）取一个，32 个线程落在 32 条不
同的段上：

\begin{lstlisting}
连续：32 线程 x 4 B = 128 B，一笔事务，利用率 100%
跨步 128 B：32 笔事务，每笔只有 4 B 有用
  带宽利用率 = 4/128 = 3.125% ≈ 1/32
  有效带宽 = 500 GB/s x 1/32 ≈ 15.6 GB/s
\end{lstlisting}

把两种地址模式画成 warp 视角的对照图，“拼成一
笔”与“散成 32 笔”的差别一眼可见（图中各画 4
条道代表 32 条）。

\begin{pdfdiagram}
  % 上：合并 —— 32 条道的地址落在同一条段内
  \node[hw label, anchor=west] at (0.6,3.62) {合并：32 条道地址连续};
  \node[hw compute, minimum width=0.6cm, minimum height=0.5cm, inner sep=1pt] (t0) at (1.9,2.95) {t0};
  \node[hw compute, minimum width=0.6cm, minimum height=0.5cm, inner sep=1pt] (t1) at (4.3,2.95) {t1};
  \node[hw compute, minimum width=0.6cm, minimum height=0.5cm, inner sep=1pt] (t2) at (6.7,2.95) {t2};
  \node[hw compute, minimum width=0.6cm, minimum height=0.5cm, inner sep=1pt] (t31) at (9.4,2.95) {t31};
  \node[hw label] at (8.05,2.95) {$\cdots$};
  \node[hw memory, minimum width=9.6cm, minimum height=0.85cm] (seg) at (5.5,1.75) {一条 128 B 段（cache 行）：$32\times4$ B 恰好填满};
  \draw[hw bus] (t0.south) -- (t0.south |- seg.north);
  \draw[hw bus] (t1.south) -- (t1.south |- seg.north);
  \draw[hw bus] (t2.south) -- (t2.south |- seg.north);
  \draw[hw bus] (t31.south) -- (t31.south |- seg.north);
  \node[hw label] at (5.5,1.05) {落在同一条段内 $\Rightarrow$ 1 笔事务，带宽利用率 100\%};
  % 下：跨步 —— 每条道各碰一条段
  \node[hw label, anchor=west] at (0.6,0.15) {跨步 128 B：每条道各碰一条段};
  \node[hw compute, minimum width=0.6cm, minimum height=0.5cm, inner sep=1pt] (s0) at (1.9,-0.55) {t0};
  \node[hw compute, minimum width=0.6cm, minimum height=0.5cm, inner sep=1pt] (s1) at (4.3,-0.55) {t1};
  \node[hw compute, minimum width=0.6cm, minimum height=0.5cm, inner sep=1pt] (s2) at (6.7,-0.55) {t2};
  \node[hw compute, minimum width=0.6cm, minimum height=0.5cm, inner sep=1pt] (s31) at (9.4,-0.55) {t31};
  \node[hw label] at (8.05,-0.55) {$\cdots$};
  \node[hw memory, minimum width=1.9cm, minimum height=0.85cm] (g0) at (1.9,-1.75) {段 0\\只用 4 B};
  \node[hw memory, minimum width=1.9cm, minimum height=0.85cm] (g1) at (4.3,-1.75) {段 1\\只用 4 B};
  \node[hw memory, minimum width=1.9cm, minimum height=0.85cm] (g2) at (6.7,-1.75) {段 2\\只用 4 B};
  \node[hw memory, minimum width=1.9cm, minimum height=0.85cm] (g31) at (9.4,-1.75) {段 31\\只用 4 B};
  \node[hw label] at (8.05,-1.75) {$\cdots$};
  \draw[hw bus] (s0.south) -- (g0.north);
  \draw[hw bus] (s1.south) -- (g1.north);
  \draw[hw bus] (s2.south) -- (g2.north);
  \draw[hw bus] (s31.south) -- (g31.north);
  \node[hw label] at (5.5,-2.6) {32 笔事务：每笔取回 128 B 只用 4 B $\Rightarrow$ 有效带宽 $\approx1/32$};
\end{pdfdiagram}
\diagramnote{合并访存与跨步访存的对照（各画 4 条道示意 32 条）。地址连续时，warp 的 32 个 4 B 请求恰好拼满一条 128 B 段，一笔事务人人有份；跨步 128 B 时，每条道落在不同的段上，32 笔事务每笔只取走 4 B 有用数据，有效带宽只剩 $1/32$。同一台显存，地址模式差出一个数量级还多。}

15.6 GB/s，比 CPU 的主存还慢一截。三种典型模式
并排看更清楚：

\begin{longtable}{L{0.26\linewidth}L{0.2\linewidth}L{0.4\linewidth}}
\toprule
warp 内地址模式 & 每指令事务数 & 有效带宽（500 GB/s 峰值） \\
\midrule
连续，32 线程 x 4 B & 1 笔 128 B & 约 500 GB/s，全额 \\\tablerowrule
跨步 16 B & 4 笔 & 约 125 GB/s，四分之一 \\\tablerowrule
跨步 128 B & 32 笔 & 约 15.6 GB/s，三十二分之一 \\
\bottomrule
\end{longtable}

这就是 GPU 版的空间局部性。存储与设备事务相关章节讲 cache 行与
突发传输时立过同一条规矩：按行取回，就用满整行。
GPU 把这条规矩放大到 warp 粒度，惩罚也放大到 32
倍——因为它更靠带宽吃饭。由此可以直接得到工程推论：结
构体数组（AoS）按字段访问天然跨步，改成数组结
构（SoA）让同名字段连续存放，往往一处改动换回
一个数量级的有效带宽；写操作同理合并，读写混合
的核心里，每个 warp 的地址模式都要过一遍这杆秤。
第三节算显存总账时，还会回到这条规矩。

\topic{推演：一次向量加法在 SIMT 下怎样跑}
把本节全部机制装进一个最小例子走一遍：
$N=2^{20}=1,048,576$ 个 float 的向量加法
\code{c[i]=a[i]+b[i]}。内核本身只有一行有效运算：
每个线程读 \code{a[i]}、读 \code{b[i]}、相加、写
\code{c[i]}。启动配置取每块 256 线程，逐层映射
到硬件。

第一步，网格与块。总块数
$1,048,576/256=4,096$，每块 $256/32=8$ 组 warp。
网格只描述工作量，不指定在哪跑；$N$ 在这里恰好
整除，不能整除时最后一组 warp 靠掩码收尾，与第
二节分歧是同一套机制。

第二步，块到 SM。块调度器把 4,096 块轮流派给 128
个 SM；每个 SM 按寄存器与共享内存的余量决定驻留
几块，假设满配 8 块，正好 $8\times8=64$ 组 warp，
占用率 100\%。4,096 块不会同时在片上，跑完一块
补一块——这个队列行为和存储与外设队列案例的命令队列是同
一张脸。

第三步，warp 到调度槽。SM 内 4 个调度器各管 16
组 warp，每拍各挑一组就绪的发射一条指令。驻留
的 64 组里，大多数在等访存回来，少数就绪——400
拍的空窗由它们轮转填满，这正是占用率演算里
“4 组起步、8 组才稳”的现场。

第四步，线程到算术道。一条 warp 指令发射时，道
$i$ 执行组内第 $i$ 个线程；线程号经
$\text{块号}\times256+\text{块内号}$ 换算成全局
下标，决定读写哪个元素。道与线程的绑定只在那一
拍有效，下一拍发射槽换上另一组 warp 的 32 个线
程。

四步映射画成一张对照图：左边是启动配置描述的软件层
级，右边是承接它的硬件——

\begin{pdfdiagram}
  % 第一步：网格 → 块调度器
  \node[hw state, minimum width=4.2cm] (g) at (2.3,0) {网格 grid：4,096 块\\（$2^{20}$ 线程 ÷ 每块 256）};
  \node[hw compute, minimum width=5.8cm] (sm) at (8.1,0) {块调度器：轮流派给 128 个 SM\\4,096 块不同时候在片上，跑完一块补一块};
  \draw[hw bus] (g.east) -- (sm.west);
  \node[hw label, anchor=south] at (4.8,0.52) {第一步};
  % 第二步：块 → SM 驻留
  \node[hw state, minimum width=4.2cm] (b) at (2.3,-1.35) {块 block：256 线程};
  \node[hw compute, minimum width=5.8cm] (res) at (8.1,-1.35) {SM 驻留：按寄存器与共享内存余量\\满配 8 块 $=$ 64 组 warp，占用率 100\%};
  \draw[hw bus] (b.east) -- (res.west);
  \node[hw label, anchor=south] at (4.8,-0.83) {第二步};
  % 第三步：warp → 调度槽
  \node[hw state, minimum width=4.2cm] (w) at (2.3,-2.7) {warp：每块 $256/32=8$ 组};
  \node[hw compute, minimum width=5.8cm] (sch) at (8.1,-2.7) {4 个调度器各管 16 组 warp\\每拍各挑一组就绪的发射一条指令};
  \draw[hw bus] (w.east) -- (sch.west);
  \node[hw label, anchor=south] at (4.8,-2.18) {第三步};
  % 第四步：线程 → 算术道
  \node[hw state, minimum width=4.2cm] (t) at (2.3,-4.05) {线程：块号 $\times256+$ 块内号\\换算成全局下标};
  \node[hw compute, minimum width=5.8cm] (lane) at (8.1,-4.05) {32 条算术道：道 $i$ 执行组内线程 $i$\\绑定只在这一拍有效};
  \draw[hw bus] (t.east) -- (lane.west);
  \node[hw label, anchor=south] at (4.8,-3.53) {第四步};
\end{pdfdiagram}
\diagramnote{一次向量加法启动的四级映射（$N=2^{20}$、每块 256 线程）。网格只描述工作量：4,096 块由块调度器轮流派给 128 个 SM；每个 SM 按资源余量驻留若干块，满配 8 块合 64 组 warp；SM 内 4 个调度器分管这些 warp，每拍各挑一组就绪的发射；一条 warp 指令发射时，道 $i$ 执行组内第 $i$ 个线程，道与线程的绑定只在这一拍有效。}

第五步，访存合并。每个 warp 读 \code{a[]} 的 32
个下标天然连续，$32\times4=128$ 字节一笔事务，
上一题的结论直接兑现；\code{b[]} 与写回
\code{c[]} 同理。若下标换算里掺了跨步，带宽立刻
按上一题的表打折扣。

最后算两笔账，看这件事总共花多少时间。指令账：

\begin{lstlisting}
每线程 2 load + 1 add + 1 store，共 4 类指令
warp 指令总数 = 1,048,576 / 32 x 4 = 131,072 条
全卡发射能力 = 128 SM x 4 调度器 = 512 条/拍
发完指令流 ≈ 131,072 / 512 = 256 拍 ≈ 0.17 微秒
\end{lstlisting}

数据账：

\begin{lstlisting}
数据量 = 读 8 MiB + 写 4 MiB = 12 MiB ≈ 12.6 MB
1 TB/s 带宽下：12.6 MB / 1 TB/s ≈ 12.6 微秒
\end{lstlisting}

0.17 微秒对 12.6 微秒，指令发射比数据到达快了近
两个数量级——这是一个纯带宽题。用“冒险、异常与性能”一章
roofline 的话说：算术强度是 $1/12\approx0.08$ 次
运算每字节，机器拐点是
$49\ \text{TFLOPS}/1\ \text{TB/s}\approx49$ 次每
字节，$0.08\ll49$，深深位于带宽区。这类内核的
优化方向因此只有一个：访存模式与有效带宽，而不
是算术。与存储与外设队列案例中的 NVMe 一次读的 trace 对照着看
更有意思：两边都是“队列下发、硬件自取、完成通
知”的事务，区别在规模——SM 内部这十几微秒里
的几百次 warp 切换对 CPU 完全不可见，CPU 只在
fence 处看到“完成”二字（第四节展开）。

\section{三、显存系统：VRAM、HBM 与数据进出}

显存是 GPU 版图里最显眼的分界线。CPU 的主存是所有设备
共享的主场，GPU 却另立一座仓库：数据先要搬进显存，才能
被上万个线程高速消费。

\topic{VRAM 身兼三职}
VRAM（显存）的第一职是帧缓冲：屏幕每个像素的颜色值躺在
显存一块连续区域里，等显示控制器按固定节拍逐行取走——
那是第四节的题目。一块 $1920\times1080$、每像素 4 字节
的帧缓冲约 8.3 MB，双缓冲再翻一倍，对今天的显存容量只
是零头。第二职是图形资源的仓库：纹理贴图、顶点缓冲、着
色器编译后的指令。这一职的访问强度远大于帧缓冲——同一
小片纹理在一帧里被成千上万个像素反复采样，命中频次以亿
计。第三职是通用计算的数据区：CUDA 程序里的数组、矩阵、
中间结果都分配在显存里，内核启动后线程直接读写，计算密
集时每个字节每毫秒都会被访问好几次。

\begin{longtable}{L{0.14\linewidth}L{0.34\linewidth}L{0.36\linewidth}}
\toprule
用途 & 装的什么 & 访问特征 \\
\midrule
帧缓冲 & 待显示的像素颜色 & 显示控制器逐行顺序读，渲染引擎随机写 \\\tablerowrule
纹理/资源 & 贴图、顶点、着色器指令 & 读多写少，同一片数据被大量像素重复读 \\\tablerowrule
通用数据 & 数组、矩阵、中间结果 & 由内核决定，读写都可能极密集 \\
\bottomrule
\end{longtable}

三职共享同一块显存，也共享同一份带宽：显示扫描每帧固定
取走约 0.5 GB/s（第四节会算），只是零头；纹理采样与通
用计算才是真正的带宽大户，两者靠片上 cache 与调度错开
峰值。显存容量也跟着三职一路膨胀：早期显卡几百 MB 就够
放帧缓冲与纹理，今天一张计算卡配几十 GB，装的是模型参
数与中间张量——第三职成了容量的主要推手。

三职的访问局部性也各不相同，这解释了 GPU 存储设计的取
舍。帧缓冲按行顺序被扫、按块随机被写，局部性最好；纹理
采样有强烈的二维邻近性——相邻像素采相邻纹素——所以
GPU 给纹理配了专用 cache 与专门的二维分块布局，把线性
地址打乱成小块，让一次采样带回来的整行数据下一拍还能用
；通用计算的局部性完全由程序员决定，写得好时比纹理还整
齐，写得差时全是指针追逐，硬件帮不上忙。同一块显存上三
种访问模式并存，是 GPU 存储子系统比 CPU 更难设计的第一
个原因；第二个原因前面已经说过——它还要大一个数量级的
带宽。

三职合在一处，“显存不够怎么办”就成了必答题。操作系统
与驱动给的答案是换页：把暂时不用的部分挪回主存，用到时
再搬回来，程序逻辑上不必修改。能跑，但性能是断崖。量级
一摆就清楚：本节后面要算的 HBM 显存带宽约 820 GB/s，
跨 PCIe 回主存只有约 30 GB/s，差 $25$ 倍以上；时延同样
差一个数量级，显存内部一次访问约几百纳秒，跨总线到主存
以微秒计。更麻烦的是触发方式：换页由缺页异常驱动，内核
运行到一半发现数据不在显存，整个网格停下来等搬运，一次
缺页就把十万线程并行攒下的吞吐优势赔进去。工程结论因此
很硬：显存容量是 GPU 程序的第一约束；放不下的数据集必
须先分块，让每一块在显存里被反复用够再换下一块，而不是
依赖换页机制补偿容量不足。第四节讲命令路径时还会看到，连“什么时候
搬”都是显式安排的，而不是等异常发生。

\topic{GDDR 与 HBM 的带宽来源}
显存带宽不是某项聪明设计的产物，而是两个朴素物理量的乘
积：带宽（B/s）$=$ 总线位数 $\times$ 每引脚传输率
$\div8$。GDDR 与 HBM 的分歧，在于各自把乘积中的哪一个
因子推到极致。GDDR 走“高频点对点”：每颗显存颗粒以
32 位接口直连 GPU，走线短、负载轻，单引脚跑到
16--24 Gbps——GDDR6 约 16 Gbps，GDDR6X 改用四电平
PAM4 编码，用两个比特压进一次电平跳变，单引脚推到
24 Gbps。十二颗颗粒围 GPU 一圈，拼出 384 位接口。这条
路线的天花板在封装边缘：384 位意味着芯片四周引出近四百
组数据线加配套的电源与地，位宽再想翻倍，封装周长不够用
了；20 Gbps 以上的信号完整性也让每一代提速都更艰难。
HBM 走另一条路：“超宽低频堆叠”。DRAM 芯片垂直堆成一
塔，用硅穿孔（TSV）上下连通，整塔与 GPU 并排安放在同一
块硅中介层上；单塔接口宽达 1024 位——是 GDDR 单颗的
32 倍——每引脚反而只要 1--2 Gbps。宽度是堆出来的：
一塔 1024 位，四塔并排就是 4096 位。

\begin{lstlisting}
GDDR6：384 位 x 16 Gbps/引脚 / 8 = 768 GB/s
GDDR6X：384 位 x 24 Gbps/引脚 / 8 = 1152 GB/s
HBM2：1024 位/塔 x 4 塔 = 4096 位
      4096 位 x 1.6 Gbps/引脚 / 8 = 819.2 GB/s，约 820 GB/s
对照：双通道 DDR5-6400 主存 = 128 位 x 6.4 Gbps / 8 ≈ 102 GB/s
\end{lstlisting}

两条路线的结构并排画出来，“哪个因子被推到极致”一
眼可辨：

\begin{pdfdiagram}
  % 左：GDDR6 高频点对点
  \node[hw label] at (2.85,0.72) {GDDR6：高频点对点};
  \node[hw memory, minimum width=1.05cm, minimum height=0.7cm, inner sep=1pt] (m0) at (0.95,0) {颗粒 0\\32 位};
  \node[hw memory, minimum width=1.05cm, minimum height=0.7cm, inner sep=1pt] (m1) at (2.1,0) {颗粒 1\\32 位};
  \node[hw label] at (2.85,0) {$\cdots$};
  \node[hw memory, minimum width=1.05cm, minimum height=0.7cm, inner sep=1pt] (m10) at (3.6,0) {颗粒\\$\cdots$};
  \node[hw memory, minimum width=1.05cm, minimum height=0.7cm, inner sep=1pt] (m11) at (4.75,0) {颗粒 11\\32 位};
  \node[hw compute, minimum width=4.9cm, minimum height=0.9cm] (gpu1) at (2.85,-1.35) {GPU：384 位接口};
  \draw[hw bus] (m0.south) -- (m0.south |- gpu1.north);
  \draw[hw bus] (m1.south) -- (m1.south |- gpu1.north);
  \draw[hw bus] (m10.south) -- (m10.south |- gpu1.north);
  \draw[hw bus] (m11.south) -- (m11.south |- gpu1.north);
  \node[hw label] at (2.85,-2.4) {12 颗 $\times$ 32 位 $=$ 384 位，$\times$ 16 Gbps $\div8=$ 768 GB/s};
  % 右：HBM2 超宽低频堆叠
  \node[hw label] at (9.0,0.72) {HBM2：超宽低频堆叠};
  \node[hw compute, minimum width=2.0cm, minimum height=1.0cm] (gpu2) at (7.3,-1.35) {GPU};
  \node[hw memory, minimum width=1.5cm, minimum height=0.72cm, inner sep=1pt] (h0) at (10.15,0) {HBM 塔\\1024 位};
  \node[hw memory, minimum width=1.5cm, minimum height=0.72cm, inner sep=1pt] (h1) at (10.15,-0.9) {HBM 塔\\1024 位};
  \node[hw memory, minimum width=1.5cm, minimum height=0.72cm, inner sep=1pt] (h2) at (10.15,-1.8) {HBM 塔\\1024 位};
  \node[hw memory, minimum width=1.5cm, minimum height=0.72cm, inner sep=1pt] (h3) at (10.15,-2.7) {HBM 塔\\1024 位};
  \draw[hw bus] (gpu2.east) -- (9.05,-1.35);
  \draw[BookMuted, line width=0.62pt] (9.05,0) -- (9.05,-2.7);
  \draw[hw bus] (9.05,0) -- (h0.west);
  \draw[hw bus] (9.05,-0.9) -- (h1.west);
  \draw[hw bus] (9.05,-1.8) -- (h2.west);
  \draw[hw bus] (9.05,-2.7) -- (h3.west);
  \node[hw label] at (9.0,-3.6) {4 塔 $\times$ 1024 位 $=$ 4096 位，$\times$ 1.6 Gbps $\div8\approx$ 820 GB/s};
  % 公式
  \node[hw label] at (5.75,-4.4) {带宽（B/s）$=$ 总线位数 $\times$ 每引脚传输率 $\div8$：GDDR 推速率到极致，HBM 把位宽堆到极致};
\end{pdfdiagram}
\diagramnote{GDDR 与 HBM 求带宽的两条路线（各画部分颗粒与塔示意）。GDDR 走高频点对点：每颗颗粒 32 位直连 GPU，单引脚 16 Gbps，十二颗拼出 384 位、768 GB/s，位宽再涨受封装周长限制。HBM 走超宽低频：DRAM 用硅穿孔堆成塔、与 GPU 同放一块硅中介层，单塔 1024 位、每引脚只要 1.6 Gbps，四塔 4096 位约 820 GB/s，宽度是用堆叠与封装成本换的。同一道公式，一个把速率推到极致，一个把位宽堆到极致。}

\begin{longtable}{L{0.13\linewidth}L{0.18\linewidth}L{0.22\linewidth}L{0.33\linewidth}}
\toprule
路线 & 每引脚速率 & 总位宽 & 带宽的物理出处 \\
\midrule
GDDR6/6X & 16--24 Gbps & 256--384 位 & 封装边缘高频点对点，位宽受引脚数限制 \\\tablerowrule
HBM2 & 1--2 Gbps & 1024 位/塔，可堆 4--8 塔 & 硅中介层超宽接口，宽度用堆叠换 \\\tablerowrule
DDR 主存 & 3--6.4 Gbps & 64 位/通道 & 双通道 128 位，带宽约 100 GB/s \\
\bottomrule
\end{longtable}

三组数字并排，结论一目了然：HBM2 的 820 GB/s 约是双通
道 DDR5 主存的 8 倍，而这 8 倍没有任何魔法——位宽之比
$4096/128=32$，引脚速率之比约 $1/4$，两者相乘正好是 8。
带宽是物理堆出来的：要宽度，就得出面积、出硅中介层、出
封装成本；要速率，就得在信号完整性上持续投入，还要承受
高频带来的功耗。HBM 每引脚速率低附带一个红利：速率低则
每比特搬运的能耗低，同样带宽下 HBM 的显存功耗明显小于
GDDR——这正是功耗顶格的计算卡偏爱 HBM、成本敏感的消
费卡坚持用 GDDR 的原因。选哪条路线是预算题，不是高下题
。HBM 的短板也要说全：硅中介层与堆叠封装把成本抬高一截
，单塔容量受堆叠层数封顶，所以 HBM 卡贵且容量档位少；
GDDR 颗粒便宜、扩容容易，代价是带宽与能效都停在更低的
档位上。

两条路线各自的演进也沿着同一道公式走。HBM 从 HBM2 的
1.6--2 Gbps 每引脚提到 HBM3 的 6.4 Gbps，位宽一塔没变，
带宽直接增到四倍，单塔从约 205 GB/s 涨到 819 GB/s——速率
这条轴上它一直在补课。GDDR 守不住位宽就守速率：GDDR6X
引入 PAM4 之后，再往下走要靠更激进的编码与更短更平的走
线，每代提速换来的 Gbps 越来越少。把两家的路线图叠在一
起看，结论还是那个：位宽必须用面积与封装换，速率必须向
信号完整性换，带宽没有第三种来源。

\topic{数据进出：PCIe DMA 与 pinned memory}
独立显存的另一面，是数据得先跨过 PCIe 才能进 GPU。搬运
不靠 CPU 逐字拷贝——那会把一个核完全占死——而是由
GPU 板上的 DMA 引擎完成：驱动把源地址、目的地址、长度写
成描述符，引擎自己读主存、写显存，完成后写状态、发中断
。这正是存储与设备事务相关章节描述符 DMA 的那张图，只是执行体从磁盘控
制器换成了 GPU。瓶颈在链路本身：PCIe 4.0 x16 每通道
16 GT/s，共 16 通道，128b/130b 编码，理论带宽
$16\times16\times128/130\div8\approx31.5$ GB/s，扣除
TLP 包开销实得约 25--30 GB/s；上一代的 3.0 x16 则只有
约 16 GB/s。对照上一题的 820 GB/s，进出 GPU 的路比显存
内部窄了约 30 倍——显存系统真正的窄门不在显存，在总
线。

有一个使用细节必须先交代：普通可分页内存不能被 DMA 直
接使用，它的物理页随时可能被操作系统换走，DMA 搬到一半
页面易主，就是数据损坏。所以可分页数据要走两趟：CPU 先
把它拷进一块锁页的中转缓冲，再由 DMA 从缓冲搬走。锁页
内存（pinned memory，页被锁定不许换出）则省掉中转，DMA
一步直达。搬 1 GiB 数据，两条路径各花多少时间，可以直
接算：

\begin{lstlisting}
可分页路径：CPU 拷 1 GiB 到中转缓冲（单流约 10 GB/s，约 107 ms）
           + DMA 搬 1 GiB（约 25 GB/s，约 43 ms）≈ 150 ms
锁页路径：  DMA 一步直达，1 GiB / 25 GB/s ≈ 43 ms
对照：     显存内部拷 1 GiB = 读 1 GiB + 写 1 GiB
           2.15 GB / 820 GB/s ≈ 2.6 ms
\end{lstlisting}

两条路径并排画出来，150 ms 与 43 ms 的差距就是两趟与一趟的差距：

\begin{pdfdiagram}
  % 上：可分页，两趟
  \node[hw label, anchor=east] at (-0.15,2.3) {可分页};
  \node[hw memory, minimum width=1.75cm, minimum height=0.95cm] (m1) at (0.9,2.3) {主存\\数据};
  \node[hw compute, minimum width=1.75cm, minimum height=0.95cm] (c1) at (3.2,2.3) {CPU 拷贝\\107 ms};
  \node[hw state, minimum width=1.75cm, minimum height=0.95cm] (b1) at (5.5,2.3) {锁页中转\\缓冲};
  \node[hw control, minimum width=1.75cm, minimum height=0.95cm] (d1) at (7.8,2.3) {DMA 引擎\\43 ms};
  \node[hw memory, minimum width=1.75cm, minimum height=0.95cm] (g1) at (10.1,2.3) {显存};
  \draw[hw bus] (m1.east) -- (c1.west);
  \draw[hw bus] (c1.east) -- (b1.west);
  \draw[hw bus] (b1.east) -- (d1.west);
  \draw[hw bus] (d1.east) -- (g1.west);
  \node[hw label, anchor=north] at (5.5,1.58) {两趟合计约 150 ms：一次整机参与的 CPU 拷贝加一趟 DMA};
  % 下：锁页直达，一趟
  \node[hw label, anchor=east] at (-0.15,0.2) {锁页};
  \node[hw memory, minimum width=1.75cm, minimum height=0.95cm] (m2) at (0.9,0.2) {主存\\（锁页）};
  \node[hw control, minimum width=1.75cm, minimum height=0.95cm] (d2) at (5.5,0.2) {DMA 引擎};
  \node[hw memory, minimum width=1.75cm, minimum height=0.95cm] (g2) at (10.1,0.2) {显存};
  \draw[hw bus] (m2.east) -- (d2.west);
  \draw[hw bus] (d2.east) -- (g2.west);
  \node[hw label, anchor=north] at (5.5,-0.52) {一趟直达约 43 ms；对照：显存内部拷 1 GiB 仅约 2.6 ms};
\end{pdfdiagram}
\diagramnote{搬 1 GiB 进显存的两条路径。可分页内存的物理页随时可能被换出，DMA 不能直接使用，必须先由 CPU 拷进锁页中转缓冲（约 107 ms）再交给 DMA（约 43 ms），两趟合计约 150 ms；锁页内存省掉中转，DMA 一步直达约 43 ms。对照线：同样的 1 GiB 在显存内部拷贝只需约 2.6 ms——跨总线比显存内部贵一个数量级以上。}

三组数字的工程含义一层比一层硬。第一层：跨总线搬运比显
存内部搬运贵一个数量级以上，43 ms 在 60 Hz 下接近三帧
——渲染循环若每帧都要从主存拉 1 GiB 新数据，帧率上限
立刻被锁在 20 fps 出头，与算力完全无关。第二层：锁页把
两趟并成一趟，省掉的是一次整机参与的 CPU 拷贝，所以大
流量传输的第一条规矩就是缓冲一律锁页。第三层：锁页的页
退出主存调度，锁页过多会严重占用整机内存资源，惯例是开
几块锁页缓冲循环复用，而不是把全部数据区都锁死。还有第
四层：小拷贝比大拷贝更糟，每次 DMA 提交有微秒级的固定
开销，若按 4 KB 一批拆着传 1 GiB，光固定开销就叠到一秒
量级（每次数微秒）——传输必须攒大块，这与“冒险、异常与性能”一章“批量摊薄固定成本”
是同一条算术。在此之上再叠一层调度：DMA 引擎与计算引擎
各自独立，拷贝队列与计算队列并行推进，第 $k$ 块数据在
总线上搬运时第 $k-1$ 块正在算，分块加异步，把 30 倍差
距藏进流水。

链路本身也在演进，但节奏值得注意：PCIe 5.0 x16 把每通
道速率翻倍到 32 GT/s，理论带宽约 63 GB/s，6.0 再翻一倍
——每代翻一倍，听起来很快；可与显存侧对比，HBM 单塔约六年增到四倍，总线始终在追赶，相对差距没有缩小。所以“
数据进了显存就别出来”这条纪律，每一代硬件都在加强而不
是放松。工程上的对应做法也越来越统一：训练框架把数据集
常驻显存，推理框架把权重锁在显存里，跨卡交换走 GPU 之
间的直连链路而不是绕回主存——都是同一个动机：PCIe 的
30 倍差距，绕开一次是一次。

最后一层便利来自统一虚拟寻址（UVA，unified virtual
addressing）：主存锁页区与显存被映射进同一张 64 位虚拟
地址空间，指针的数值本身就标明了数据在哪一端，驱动按地
址区间路由访问。程序员的实惠是拷贝函数不必分别声明“主
存指针”“显存指针”，同一个指针两端通用；硬件的实惠是
地址翻译表只有一套。但务必记住，UVA 统一的是编址，不是
速度：指针两端通用，带宽两端照旧差 30 倍。

把整道题收拢成操作清单，只有三条。其一，算得久的数据一
次传入、常驻显存，传输成本按使用次数摊薄；其二，必须频
繁进出的数据走锁页缓冲加异步 DMA，与计算重叠，把总线时
间藏进显存时间里；其三，放不下的数据按块调度，块的大小
至少要让“算这块的时间”显著大于“搬这块的时间”——搬
算比不足 $1:10$ 的分块方案，等于花钱买总线。这三条与算
力无关，却决定了多数真实程序的实际速度。

\topic{共享内存与寄存器：GPU 的两级近端存储}
显存带宽再大，离算术单元仍嫌远：一次显存往返约四到八百拍，
操作数若每个都现取现用，再宽的带宽也被时延拖
空。GPU 的答案是把两级近端存储摆得极近、给得极多。第一
级是寄存器堆：每个 SM 配 64K 个 32 位寄存器、共 256 KB，
按线程静态瓜分——驻留 2048 个线程时，每线程分到 32
个；驻留 1024 个时，每线程 64 个。寄存器由编译器直接编
进指令，访问不占额外周期，是真正意义上的零等待。第二级
是共享内存（shared memory）：每个 SM 数十到两百 KB 的片
上 SRAM，由同一线程块内的线程共同寻址，时延约二三十个
周期，介于寄存器与显存之间。与这两级形成对照的是
L1/L2 cache：它们存在、也有用，但由硬件自动管理，容量
与命中率都不由程序员直接安排。

\begin{longtable}{L{0.15\linewidth}L{0.22\linewidth}L{0.18\linewidth}L{0.3\linewidth}}
\toprule
层级 & 容量量级 & 时延量级 & 谁在管理 \\
\midrule
寄存器 & 每 SM 256 KB & 零额外周期 & 编译器静态分配，按线程私有 \\\tablerowrule
共享内存 & 每 SM 数十至两百 KB & 约 20--30 周期 & 程序员显式读写，块内共享 \\\tablerowrule
L1/L2 cache & 每 SM 百余 KB / 整卡数十 MB & 约 30 / 200 周期 & 硬件自动，与 CPU cache 同理 \\\tablerowrule
显存 VRAM & 数 GB 至数十 GB & 约四到八百拍 & 分配器管理空间，带宽见前文 \\\tablerowrule
主存（跨 PCIe） & 随整机 & 微秒级 & 操作系统与驱动，换页触发 \\
\bottomrule
\end{longtable}

五个层级连同管理者排成一列，GPU 与 CPU 的差别全在
右列：

\begin{pdfdiagram}
  \node[hw label] at (2.9,0.78) {层级（向下容量更大、时延更长）};
  \node[hw label] at (8.6,0.78) {谁在管理};
  \node[hw memory, minimum width=3.4cm, minimum height=0.8cm] (l1) at (2.2,0) {寄存器：每 SM 256 KB\\零额外周期};
  \node[hw memory, minimum width=4.0cm, minimum height=0.8cm] (l2) at (2.5,-0.95) {共享内存：每 SM 数十至两百 KB\\约 20--30 周期};
  \node[hw memory, minimum width=4.6cm, minimum height=0.8cm] (l3) at (2.8,-1.9) {L1/L2 cache：百余 KB / 数十 MB\\约 30 / 200 周期};
  \node[hw memory, minimum width=5.2cm, minimum height=0.8cm] (l4) at (3.1,-2.85) {显存 VRAM：数 GB 至数十 GB\\约 400--800 周期};
  \node[hw memory, minimum width=5.8cm, minimum height=0.8cm] (l5) at (3.4,-3.8) {主存（跨 PCIe）：随整机\\微秒级};
  \draw[BookRule, dotted, line width=0.5pt] (l1.east) -- (7.1,0);
  \draw[BookRule, dotted, line width=0.5pt] (l2.east) -- (7.1,-0.95);
  \draw[BookRule, dotted, line width=0.5pt] (l3.east) -- (7.1,-1.9);
  \draw[BookRule, dotted, line width=0.5pt] (l4.east) -- (7.1,-2.85);
  \draw[BookRule, dotted, line width=0.5pt] (l5.east) -- (7.1,-3.8);
  \node[hw label, anchor=west] at (7.3,0) {编译器静态分配，按线程私有};
  \node[hw label, anchor=west] at (7.3,-0.95) {程序员显式读写，块内共享};
  \node[hw label, anchor=west] at (7.3,-1.9) {硬件自动，与 CPU cache 同理};
  \node[hw label, anchor=west] at (7.3,-2.85) {分配器管理空间};
  \node[hw label, anchor=west] at (7.3,-3.8) {操作系统与驱动，换页触发};
  \node[hw label] at (5.5,-4.7) {容量最大、收益最高的两级近端存储均需要软件显式安排——GPU 没有硬件自动管理机制};
\end{pdfdiagram}
\diagramnote{GPU 存储层级与各级管理者。寄存器与共享内存这两级近端存储由编译器与程序员显式安排：寄存器按线程静态瓜分，多用一个变量就可能挤掉一整组驻留线程；共享内存由程序员写代码搬运与复用。L1/L2 仍由硬件自动管理，再往外是显存与跨 PCIe 的主存。与 CPU 全由硬件 cache 自动管理不同，GPU 将近端两级交由程序员显式管理——安排合理时效率更高，安排不当时也没有自动补偿机制。}

与 CPU 的存储层次对照，最明显的差别在“谁在管理”那一
列。CPU 的 cache 对程序员透明，访问模式规整的循环可由硬
件 cache 自动利用局部性；GPU 容量最大、收益最高的两级近端存储，均需要软件显式安
排，而且各有代价。寄存器按线程配额：内核变量过多时会
超出配额，编译器把超额部分溢出到显存，性能随之下降；寄存
器用量还反过来决定驻留线程数——每线程多用一倍寄存器
，可用于隐藏时延的并发线程数就少一半，这一关系已在第二节讨论占用率时量化。共享内存更彻底：矩阵乘法先由程序员写代码把子矩
阵从显存搬进共享内存，块内线程可将这块数据复用几十至上百
次，再换下一块；搬多少、何时搬、谁读哪一块，都由程序显式指定
。可以说 GPU 将 cache 管理从硬件自动控制转为程序员显式控制
：控制权更大，安排合理时效率更高，安排不当时也没有硬件自动补偿机制。
这正是 GPU 优化文献里“分块（tiling）”出现的频率远高
于任何指令技巧的原因——下一题的演算会给出它的定量回
报。

给一组数字感受配额的刚性。一个 SM 的 64K 个寄存器，驻
留 2048 线程、每线程 32 个寄存器时恰好用满：$2048\times
32=65536$，一个不多一个不少。内核若被编译出 33 个寄存
器，驻留就掉到 $65536/33\approx1986$，按 warp 粒度向下取整
为 62 组——一个变量之差，驻留从 64 组掉到 62 组，占用率掉
了一档；若寄存器按更粗的分配粒度取整，掉得还会更多。共享内
存同样刚性：每块申请 48 KB、SM 上共有 100
KB 时，同驻的块数就是两块，第三块再怎么小也挤不进来。
这些取舍在 CPU 世界里由硬件 cache 的替换策略自动平滑，
在 GPU 世界里则全部以整数台阶的形式出现在程序员面前。

\topic{带宽预算演算：向量操作为什么是带宽瓶颈}
现在把本节全部数字收进“冒险、异常与性能”一章的 roofline。判定式只有一
条：算术强度 $I$（每字节流量配多少 FLOP）与拐点
$I^{*}=$ 峰值算力 $\div$ 带宽比较，小于拐点即带宽瓶颈
。取一张有代表性的计算卡：峰值 fp32 算力 30 TFLOP/s，
HBM 带宽 820 GB/s，拐点 $I^{*}=30000/820\approx36.6$
FLOP/B。再看两个 kernel。向量加法 \code{C[i]=A[i]+B[i]}
：每元素读 4 B、再读 4 B、写回 4 B，共 12 字节流量，只
做 1 次加法，$I=1/12\approx0.08$ FLOP/B；哪怕按最宽松
的一字节数据计，每次运算配读、读、写三次访问，也只有
3 字节流量、$I\approx0.33$，离拐点仍差两个数量级。矩阵
乘法 \code{C=A*B}：$n\times n$ 时运算 $2n^{3}$ 次、流量
最少 $12n^{2}$ 字节，$I\approx n/6$；$n=1024$ 时约
170 FLOP/B，远在拐点右侧。

\begin{lstlisting}
拐点 I* = 30 TFLOP/s / 820 GB/s ≈ 36.6 FLOP/B
向量加(fp32)：I = 1/12 ≈ 0.08 FLOP/B
  可达性能 = 820 GB/s x 1/12 ≈ 68 GFLOP/s，约为峰值的 0.23%
矩阵乘(n=1024)：I = 2n^3 / 12n^2 = n/6 ≈ 170 FLOP/B
  带宽允许 820 x 170 ≈ 140 TFLOP/s > 峰值 30，算力瓶颈，峰值用满
\end{lstlisting}

\begin{longtable}{L{0.16\linewidth}L{0.18\linewidth}L{0.22\linewidth}L{0.28\linewidth}}
\toprule
kernel & 算术强度 & 可达性能 & 瓶颈位置 \\
\midrule
向量加法 & 约 0.08 FLOP/B & 约 68 GFLOP/s & 带宽，峰值用不到 1\% \\\tablerowrule
矩阵乘法 & 约 170 FLOP/B & 30 TFLOP/s & 算力，峰值可以真正用满 \\
\bottomrule
\end{longtable}

把拐点与两个 kernel 的落点画在同一张图上：

\begin{pdfdiagram}
  % 双对数坐标：x = (log10 I + 1.5)*2.5，y = (log10 P + 1.8)*1.35
  \draw[BookMuted, line width=0.6pt, ->] (0,0) -- (10.3,0);
  \draw[BookMuted, line width=0.6pt, ->] (0,0) -- (0,4.9);
  \node[hw label, anchor=north] at (5.1,-0.62) {算术强度 $I$（FLOP/B，对数）};
  \node[hw label, rotate=90, anchor=south] at (-1.05,2.4) {可达性能（TFLOP/s，对数）};
  \foreach \x/\t in {1.25/0.1, 3.75/1, 6.25/10, 8.75/100} {
    \draw[BookMuted, line width=0.5pt] (\x,0.06) -- (\x,-0.06);
    \node[hw label, anchor=north] at (\x,-0.1) {\t};
  }
  \foreach \y/\t in {1.08/0.1, 2.43/1, 3.78/10} {
    \draw[BookMuted, line width=0.5pt] (0.06,\y) -- (-0.06,\y);
    \node[hw label, anchor=east] at (-0.1,\y) {\t};
  }
  % 屋顶线：P = 0.82 * I 上升到拐点，之后贴平 30
  \draw[BookAccent, line width=1.0pt] (0.5,0.55) -- (7.66,4.43) -- (9.9,4.43);
  % 拐点参考线
  \draw[BookMuted, line width=0.5pt, dashed] (7.66,0) -- (7.66,4.43);
  \draw[BookMuted, line width=0.5pt, dashed] (0,4.43) -- (7.66,4.43);
  \node[hw label, anchor=north] at (7.66,-0.1) {拐点 36.6};
  \node[hw label, anchor=east] at (-0.1,4.43) {30};
  % 两个落点
  \fill[BookRed] (1.01,0.85) circle (0.07);
  \node[hw label, anchor=west] at (1.25,0.48) {向量加：$I\approx0.08$，68 GFLOP/s（峰值 0.23\%）};
  \fill[BookTeal] (9.32,4.43) circle (0.07);
  \node[hw label, anchor=east] at (9.2,4.72) {矩阵乘：$I\approx170$，用满 30 TFLOP/s};
\end{pdfdiagram}
\diagramnote{双对数坐标下的 roofline：斜段是带宽屋顶（可达性能 = 带宽 $\times$ 算术强度），平段是算力屋顶 30 TFLOP/s，拐点 $I^{*}=30000/820\approx36.6$ FLOP/B。向量加的强度只有约 0.08 FLOP/B，被带宽屋顶压在 68 GFLOP/s——峰值的 0.23\%；矩阵乘约 170 FLOP/B，远在拐点右侧，峰值真正用满。横轴每往左一格，同样的算力折扣再重一倍。}

“向量操作是带宽瓶颈”至此有了定量说法：不是向量单元不
够快，而是数据供给永远跟不上——一次加法要喂 12 个字节
，喂得再快的显存也只能让算术单元空等。工程推论有三条。
其一，逐元素的简单操作（向量加、缩放、激活函数）在 GPU
上天生用不满算力，优化方向是减少趟数：把连续几个逐元素
操作融合成一个内核，中间结果留在寄存器里不写回显存，强
度立刻按省下的趟数翻倍。其二，只有高强度操作（矩阵乘、
卷积）值得为算力付费，AI 加速硬件围绕矩阵乘法展开，正
是这条算术的产业化。其三，拐点随机器移动：换 HBM3、带
宽 3 TB/s 的新卡，拐点降到约 10 FLOP/B，向量加法依旧深
在左侧，一分算力红利也拿不到——“冒险、异常与性能”一章那句“左侧
kernel 为带宽付费”，在 GPU 上比在任何平台都极端。读规
格书时也该养成对应的习惯：看到一张新卡，先查带宽与峰值
算力算出拐点，再问自己手上的 kernel 强度多少，两个数就
能判断这张卡对你有没有用。

把推论一演算到底，收益是肉眼可见的。两步连续的逐元素操
作，比如先算 \code{D[i]=A[i]+B[i]} 再算 \code{E[i]=D[i]
*2}，分两趟跑：流量是 $5\times4=20$ 字节（读 A、读 B、
写 D、读 D、写 E），做 2 次运算，$I=0.1$ FLOP/B。融合成
一个内核后，D 留在寄存器里不落显存，流量只剩 $3\times4
=12$ 字节，$I=2/12\approx0.167$——强度提高 $67\%$，可
达性能同步提高 $67\%$，而算术单元一行代码都没改。深度
学习框架里的“算子融合”做的正是这件事：把激活、归一化
、逐元素变换统统并进前一个高强度内核的尾巴上，让每一次
显存读写服务的运算尽量多。在 roofline 左侧，省字节永远
比加算力值钱——这是显存这一节最值得带走的一句话。

\section{四、命令路径：命令环、fence 与显示扫描}

GPU 与 CPU 的协作不靠函数调用，而靠一条与存储与外设队列案例中的 NVMe
神似的命令路径：驱动把命令写进队列、写一下门铃，GPU 自
己取命令、自己执行、自己报告完成。显示一侧另有一套不受
渲染指挥的固定节拍。

\topic{命令环缓冲与 doorbell}
驱动提交一批 GPU 工作的最小单元是命令包：几十到几百字
节，头部是操作码与长度，后面跟着参数——“把寄存器 X
设为 Y”“把显存 A 区拷到 B 区”“用这组参数启动一个内核
”“绘制这组三角形”。命令包不直接写给 GPU，而是按序写
进主存里的一块环缓冲（ring buffer）；写完一批，驱动做
一次 MMIO 写，把新的尾指针送进 GPU 前端的 doorbell 寄存
器。这一下就是全部通知：doorbell 不携带数据，只宣布“
环上多了东西”。GPU 一侧，命令处理器比较头尾指针，发现
差距后经 DMA 把新命令包取回片上，逐包解码、分派给对应
引擎——渲染、拷贝、计算各有引擎，也各有各的环；每消
化一段，它把读指针写回主存，驱动据此回收环上的槽位，环
满时驱动只能停下来等 GPU 消费。

这条提交—执行—回收的环路画出来，与存储与设备事务相关章节、第十三
章是同一副骨架：

\begin{pdfdiagram}
  \node[hw memory, minimum width=5.0cm, minimum height=1.0cm] (ring) at (5.5,0.4) {主存命令环：命令包 命令包 $\cdots$};
  \node[hw state, minimum width=2.4cm, minimum height=1.0cm] (drv) at (1.6,-1.7) {CPU：驱动};
  \node[hw compute, minimum width=3.0cm, minimum height=1.0cm] (gpu) at (9.8,-1.7) {GPU 前端\\命令处理器};
  % ① 命令包写入环
  \draw[hw bus] (drv.north) -- (1.6,0.4) -- (ring.west);
  \node[hw label, anchor=south] at (2.3,0.52) {① 写命令包};
  % ② doorbell：只携带新尾指针
  \draw[hw return] (drv.east) -- (gpu.west);
  \node[hw label, anchor=south] at (5.5,-1.48) {② MMIO 写 doorbell：只携带新尾指针};
  % ③ 前端按头尾指针差 DMA 取命令
  \draw[hw bus] (ring.east) -- (9.8,0.4) -- (gpu.north);
  \node[hw label, anchor=south] at (9.15,0.52) {③ DMA 取命令};
  % ④ 读指针写回，驱动回收槽位
  \draw[hw return] (gpu.south) -- (9.8,-2.85) -- (1.6,-2.85) -- (drv.south);
  \node[hw label, anchor=north] at (5.5,-2.95) {④ GPU 写回读指针，驱动回收环上槽位};
\end{pdfdiagram}
\diagramnote{GPU 命令环的提交环路。驱动把命令包按序写进主存环缓冲（①），屏障之后一次 MMIO 写把新尾指针送进 doorbell（②）——doorbell 只宣布“环上多了东西”，不携带数据；GPU 前端比较头尾指针，经 DMA 把新命令取回片上解码、分派给渲染/拷贝/计算引擎（③）；每消化一段把读指针写回主存，驱动据此回收槽位（④）。环不动、指针动，通知可合并、丢失可自愈，与 NVMe 提交队列是同一副骨架。}

把这段与存储与设备事务相关章节的描述符环并排，是同一张图的第三次出现：
环不动、指针动；doorbell 只通知、不携带数据；设备按头
尾差补读，偶尔错过一次通知也无妨。与存储与外设队列案例中的 NVMe 提交
队列的差别在细节而不在结构：NVMe 队列项定长 64 字节，
GPU 命令包变长；NVMe 常按 CPU 核绑多对队列，GPU 按上下
文与引擎分环。GPU 的特殊处只是规模：一批命令包可以上千
，覆盖一帧的全部绘制，而 NVMe 一条命令就是一次读写。所
以存储与外设队列案例“先写命令项、再写 doorbell，中间隔一道屏障”
的次序约定，在 GPU 这里逐字成立，理由也相同：通知必须
在被通知的内容可见之后发出。

为什么要环加门铃，而不是每来一条命令就叫 GPU 一次？算
一笔时延账：一次 MMIO 写加 GPU 响应的往返是微秒级，而
一条绘图命令的执行可能只有几百纳秒——逐条握手等于让
GPU 每工作一份就空等十份。环把握手摊薄：驱动攒够一批再
写一次门，门的成本被整批分摊；GPU 流水取命令，取下一批
与执行当前批重叠。工程推论：提交的开销与批次数成正比、
与命令数几乎无关，所以 GPU 编程批量越大越好——渲
染 API 鼓励“一次录一大段命令缓冲再提交”，正是这条算
术的直接应用。

把 GPU 命令环与存储与外设队列案例中的 NVMe 提交队列逐项对照，结构
上的同构就看得更清楚——学过一个，另一个的机制几乎可
以全部推导出来：

\begin{longtable}{L{0.18\linewidth}L{0.32\linewidth}L{0.32\linewidth}}
\toprule
维度 & NVMe 提交队列 & GPU 命令环 \\
\midrule
队列项 & 定长 64 字节命令 & 变长命令包，几十到几百字节 \\\tablerowrule
通知方式 & 写 SQ 尾 doorbell & 写 doorbell 寄存器，内容是新尾指针 \\\tablerowrule
取命令 & 设备经 DMA 主动取 & 前端经 DMA 主动取 \\\tablerowrule
回收依据 & CQ 完成项 + 写 CQ 头 doorbell & GPU 写回读指针，驱动回收槽位 \\\tablerowrule
队列组织 & 按 CPU 核绑多对队列 & 按上下文与引擎分环 \\
\bottomrule
\end{longtable}

还有一层二级结构值得点名：应用先把命令录进“命令缓冲”
这种可整块复用的对象，提交时再把整块挂到环上。环只是搬
运通道，命令缓冲才是复用单位——一帧里不变的绘制序列
可以录一次、提交多次，驱动省掉重新编码的开销。这与第五
章“描述符预先建好、传输时只改指针”是同一方法：把构造的
成本与提交的成本分开，贵的只做一次。

把两级结构画出来，复用单位与搬运通道各归各位：

\begin{pdfdiagram}
  \node[hw state, minimum width=1.8cm] (drv) at (1.0,-0.1) {CPU：驱动\\（应用录制）};
  \node[hw memory, minimum width=3.1cm] (bufA) at (4.15,1.15) {命令缓冲 A\\一帧不变的绘制序列};
  \node[hw memory, minimum width=3.1cm] (bufB) at (4.15,-1.35) {命令缓冲 B\\另一组可复用序列};
  \node[hw memory, minimum width=3.2cm, minimum height=1.5cm] (ring) at (7.9,-0.1) {主存命令环\\（搬运通道）\\命令包 命令包 $\cdots$};
  \node[hw compute, minimum width=1.7cm] (gpu) at (10.6,-0.1) {GPU\\前端};
  % 录一次
  \draw[hw bus] (drv.east) -- (2.35,-0.1);
  \draw[BookMuted, line width=0.62pt] (2.35,1.15) -- (2.35,-1.35);
  \draw[hw bus] (2.35,1.15) -- (bufA.west);
  \draw[hw bus] (2.35,-1.35) -- (bufB.west);
  \node[hw label, anchor=west] at (2.55,0.5) {录一次};
  % 提交＝整块挂接到环
  \draw[hw bus] (bufA.east) -- (6.5,1.15) -- (6.5,0.65);
  \draw[hw bus] (bufB.east) -- (6.5,-1.35) -- (6.5,-0.85);
  \node[hw label, anchor=west] at (6.65,1.15) {提交＝整块挂接};
  \node[hw label, anchor=west] at (6.65,-1.35) {可提交多次};
  % GPU 前端取命令
  \draw[hw bus] (ring.east) -- (gpu.west);
  \node[hw label, anchor=south] at (9.55,0.72) {DMA 取命令};
  % doorbell：只带新尾指针
  \draw[hw return] (drv.south) -- (1.0,-2.5) -- (10.6,-2.5) -- (gpu.south);
  \node[hw label, anchor=north] at (5.8,-2.6) {doorbell：只带新尾指针};
\end{pdfdiagram}
\diagramnote{命令缓冲与命令环的二级结构。应用先把命令录进命令缓冲这种可整块复用的对象（录一次），提交时再把整块挂到主存命令环上——一帧里不变的绘制序列录一次、提交多次，驱动省掉重新编码的开销。环只是搬运通道：doorbell 照样只带新尾指针，GPU 前端经 DMA 取命令。构造的成本与提交的成本分开，贵的只做一次，与存储与设备事务相关章节“描述符预先建好、传输时只改指针”是同一方法。}

doorbell 还有两个容易忽略的性质，都来自“内容只是指针”
这一件事。其一，它可以合并：驱动连续提交三批、写三次门
，GPU 看到的只是尾指针停在最后的位置，一次取走全部三批
——中间两次通知被自然吸收，不需要任何去重逻辑。其二
，它可以丢失后自愈：万一某次门铃写因为种种原因没生效，
下一次写仍会把尾指针推到位，GPU 按指针差补读，错过的那
批命令照样被取走。通知不携带状态，状态全在环与指针里，
这正是这套机制健壮与省事的共同来源。

\begin{lstlisting}
// 注释（推进）：逐行追踪发起者、所有权和可见性；请求被接受不等于事务完成。
// 注释（边界）：以采样沿、状态位或 completion 为准，再允许消费者使用结果。
驱动侧提交一批命令的序列：
1. 查环上剩余空间（尾指针与 GPU 写回的读指针之差）
2. 把命令包依次写入环缓冲（普通主存写，可被写合并）
3. 写屏障，保证命令包先于通知对 GPU 可见
4. 写 doorbell 寄存器（一次 MMIO 写，内容是新尾指针）
5. GPU 前端经 DMA 取走新命令，逐包解码、分派给引擎
6. GPU 推进读指针并写回主存，驱动回收环上槽位
\end{lstlisting}

\topic{fence 与完成语义}
doorbell 写完只代表“命令已交出”，不代表“任务已经完成”。
完成语义由 fence 提供：一批命令的末尾挂一条 fence 命令
，引擎执行到它，就把一块约定的内存位置写成预定值，可选
地再触发一次中断。CPU 或后续命令流读这个位置，就知道“
第 N 批到此全部完成”。注意 fence 的粒度：它不是每条命
令一个完成项，而是一道“到此为止”的水线。与存储与外设队列案例
NVMe 完成队列对照：NVMe 给每条命令回一项 CQ 完成项，
GPU 用一次内存写确认整批。粒度差别的理由很现实——GPU
一批命令动辄上千，逐条完成通知本身就是一笔可观的 PCIe
流量，fence 把完成语义压缩成一个字。中断也按同样思路省
：一批一中断，而不是一条一中断，否则光中断处理就把 CPU
占满。

布尔 fence 只有“未完成/已完成”两态，用完必须重置才能
复用，表达力很快见底：想表达“第 87 批完成、第 88 批还
差拷贝引擎”就捉襟见肘。时间线信号量（timeline
semaphore）把两态换成一个单调递增的 64 位计数器：每个
完成点把它推到更大的值，等待方声明“等到计数不小于 K”
。单调递增让旧值永不复现，等待条件因此没有歧义；K 可以
是一帧、一批、一个流水级，同一个对象上还能叠加任意多个
等待者与推进者。多引擎协作于是有了统一语言：渲染引擎推
到 87 之后，显示侧只等“$\geq87$”，拷贝引擎继续往 88
推，谁也不必知道别人的存在。Vulkan 与 D3D12 的现代同步
全部时间线化，布尔 fence 只留在兼容性角落里。

\begin{lstlisting}
布尔 fence：    值只有 0/1 两态；等待"为 1"；复用前须重置
时间线信号量：  值单调递增；等待"达到 K"；K = 帧号/批号/流水级
对照 NVMe：     CQ 完成项 = 每条命令回一条消息
               fence     = 一次内存写确认"到此全部完成"
\end{lstlisting}

fence 还有一个不显眼但天天在用的角色：资源回收的依据。
命令环上的槽位、这一帧用到的常量缓冲、中间纹理，都要等
对应 fence 到位才能安全复用——驱动并不“等”它，只是
定期查询，把已过水线的资源放回空闲池。完成语义在 GPU
世界里于是分成两层：给程序员的那层是“这批任务已经完成”，
给资源管理器的那层是“这块内存可以再次投入使用了”。

举一个多引擎协作的完整例子，把两种语义都用上。第 87 帧
的工作分三段：拷贝引擎先把下一帧的纹理搬进显存，完成后
把时间线信号量推到 86；渲染引擎提交时声明“等到
$\geq86$ 再开始执行”，完成后把同一个信号量推到 87；显示侧提
交翻页请求时声明“等到 $\geq87$”。三段命令分属三个引
擎的三条环，谁也不等谁地并发推进，次序全靠计数器表达；
CPU 从头到尾没有参与同步，它只是在提交第 90 帧前查了一
眼计数器，确认 87 号之前的资源都可以回收。这就是时间线
信号量比布尔 fence 灵活的根源：布尔 fence 表达一个点，
时间线表达整条进度轴。

把第 87 帧这个例子画成时间线，“谁等谁”全写在计数
器上：

\begin{waveform}[x=0.8cm,y=0.85cm]
  % 参考线：两个推进点
  \draw[BookRule, line width=0.45pt, dashed] (5,0.7) -- (5,6.3);
  \draw[BookRule, line width=0.45pt, dashed] (9,0.7) -- (9,6.3);
  \node[hw label, anchor=west] at (5.12,6.25) {拷贝完成：推到 86};
  \node[hw label, anchor=west] at (9.12,6.25) {渲染完成：推到 87};
  % 拷贝引擎
  \node[hw label, anchor=east] at (0.6,5.5) {拷贝引擎};
  \draw[BookTeal, fill=BookCodeBg, line width=0.8pt] (1,5.2) rectangle (5,5.8);
  \node at (3,5.5) {搬下一帧纹理};
  % 渲染引擎
  \node[hw label, anchor=east] at (0.6,4.15) {渲染引擎};
  \draw[BookRule, dashed, line width=0.7pt] (1,3.85) rectangle (5,4.45);
  \node[hw label] at (3,4.15) {等计数 $\geq86$};
  \draw[BookTeal, fill=BookCodeBg, line width=0.8pt] (5,3.85) rectangle (9,4.45);
  \node at (7,4.15) {渲染第 87 帧};
  % 显示侧
  \node[hw label, anchor=east] at (0.6,2.8) {显示侧};
  \draw[BookRule, dashed, line width=0.7pt] (1,2.5) rectangle (9,3.1);
  \node[hw label] at (5,2.8) {等计数 $\geq87$};
  \draw[BookTeal, fill=BookCodeBg, line width=0.8pt] (9,2.5) rectangle (12,3.1);
  \node at (10.5,2.8) {翻页};
  % 信号量计数
  \node[hw label, anchor=east] at (0.6,1.45) {信号量};
  \draw[BookTeal, fill=BookCodeBg, line width=0.8pt] (1,1.15) rectangle (5,1.75);
  \node at (3,1.45) {85};
  \draw[BookTeal, fill=BookCodeBg, line width=0.8pt] (5,1.15) rectangle (9,1.75);
  \node at (7,1.45) {86};
  \draw[BookTeal, fill=BookCodeBg, line width=0.8pt] (9,1.15) rectangle (12,1.75);
  \node at (10.5,1.45) {87};
  % 时间轴
  \draw[BookMuted, line width=0.6pt, ->] (1,0.5) -- (12.4,0.5);
  \node[hw label, anchor=west] at (12.5,0.5) {时间};
\end{waveform}
\diagramnote{时间线信号量上的多引擎协作（第 87 帧）。拷贝引擎把下一帧纹理搬进显存，完成就把单调递增的计数从 85 推到 86；渲染引擎声明“等到 $\geq86$ 再开始执行”，完成后将计数推进到 87；显示侧等 $\geq87$ 再翻页。三段命令分属三个引擎的三条环，并发推进、次序全靠计数器表达；CPU 不参与同步，只在回收资源前查一眼计数。单调递增让旧值永不复现，“等到不小于 K”没有歧义——布尔 fence 表达一个点，时间线表达整条进度轴。}

\topic{CPU-GPU 同步的三种模式}
知道“怎么查完成”之后，下一个问题是“查的时候 CPU 干
什么”。三种模式对应三种答案。轮询等待：CPU 循环读
fence 所在的内存位置，直到值到位。发现完成的时延最低
——只是一次 cache 行的读——代价是整个核被占死，还
顺带把那行 cache 踢来踢去。它适合提交后马上要结果的短
等待，比如读回一小段统计值。事件/中断：CPU 在 fence 上
登记等待后睡眠，GPU 写 fence 后触发中断唤醒。CPU 让出
来去干别的，代价是多一次中断处理与进程唤醒，时延增加数
微秒到数十微秒，适合毫秒级以上的长等待。流水化：干脆不
等——CPU 提交第 N 帧后不再碰第 N 帧的任何资源，转身开
始录第 N+1 帧的命令；fence 只在回收资源时查询，从不作
为等待点。

\begin{longtable}{L{0.14\linewidth}L{0.2\linewidth}L{0.24\linewidth}L{0.26\linewidth}}
\toprule
模式 & CPU 的占用 & 完成发现的时延 & 适用场景 \\
\midrule
轮询等待 & 整核占死 & 最低，一次 cache 读 & 微秒级短等待、读回小结果 \\\tablerowrule
事件/中断 & 让出，可干别的 & 多一次中断与唤醒 & 毫秒级以上的长等待 \\\tablerowrule
流水化 & 从不等待 & 不查询，fence 只管回收 & 渲染主循环，吞吐优先 \\
\bottomrule
\end{longtable}

流水化值得再说一层，因为它是唯一随负载变好的模式。提交
深度一帧：CPU 与 GPU 错开一整帧并行，吞吐拉满，输入到
显示的时延多出一帧上下。深度两到三帧：GPU 永远有活干，
偶发的慢帧被队列吸收，掉帧更平滑，时延再多一两帧——主
机游戏常选这里，竞技游戏宁可锁深度一。深度不是免费的：
它用内存（每级一套资源）与时延换吞吐。流水化也有一个天
敌：读回。一旦 CPU 要拿 GPU 算出的结果做下一步决策——
比如把计数器的值读回来决定后续提交多少活——就不得不
排空整条流水，攒下的深度一次赔光；所以性能敏感的引擎把
这类读回排到帧边界，或者干脆用上一帧的结果。

把深度的时延代价量化：60 Hz 下每级深度约添 16.7 ms 的输
入到显示时延，深度三帧就是约 50 ms——手快的玩家对 30
ms 以上的差距已有明确体感，这是竞技游戏把深度压到一的
全部理由。反过来，吞吐侧的账同样直接：深度为零（提交就
等）时，GPU 在每次提交边界都要空转一个提交往返的微秒级
空隙，一帧几十个边界就占用一部分帧预算。深度的选择于是
是一个明确的工程旋钮：以毫秒为单位买帧间隔的平滑，买不
买、买几毫秒，看应用对输入时延的容忍度。

把三种提交深度画成时间轴，时延与吞吐各在哪一头就清楚了：

\begin{waveform}[x=0.8cm,y=0.85cm]
  % 深度 1
  \node[hw label, anchor=east] at (0.75,7.1) {深度 1 · CPU};
  \draw[BookTeal, fill=BookCodeBg, line width=0.8pt] (1,6.85) rectangle (4,7.35);
  \node at (2.5,7.1) {录 $N$};
  \draw[BookTeal, fill=BookCodeBg, line width=0.8pt] (4,6.85) rectangle (7,7.35);
  \node at (5.5,7.1) {录 $N{+}1$};
  \draw[BookTeal, fill=BookCodeBg, line width=0.8pt] (7,6.85) rectangle (10,7.35);
  \node at (8.5,7.1) {录 $N{+}2$};
  \node[hw label, anchor=east] at (0.75,6.0) {深度 1 · GPU};
  \draw[BookAccent, fill=BookCodeBg, line width=0.8pt] (4,5.75) rectangle (7,6.25);
  \node at (5.5,6.0) {执行 $N$};
  \draw[BookAccent, fill=BookCodeBg, line width=0.8pt] (7,5.75) rectangle (10,6.25);
  \node at (8.5,6.0) {执行 $N{+}1$};
  \draw[BookAccent, fill=BookCodeBg, line width=0.8pt] (10,5.75) rectangle (13,6.25);
  \node at (11.5,6.0) {执行 $N{+}2$};
  \draw[BookMuted, line width=0.6pt, <->] (1,5.35) -- (7,5.35);
  \node[hw label, anchor=north] at (4,5.28) {输入到显示 $+1$ 帧 $\approx16.7$ ms};
  % 深度 2
  \node[hw label, anchor=east] at (0.75,4.5) {深度 2 · CPU};
  \draw[BookTeal, fill=BookCodeBg, line width=0.8pt] (1,4.25) rectangle (4,4.75);
  \node at (2.5,4.5) {录 $N$};
  \draw[BookTeal, fill=BookCodeBg, line width=0.8pt] (4,4.25) rectangle (7,4.75);
  \node at (5.5,4.5) {录 $N{+}1$};
  \draw[BookTeal, fill=BookCodeBg, line width=0.8pt] (7,4.25) rectangle (10,4.75);
  \node at (8.5,4.5) {录 $N{+}2$};
  \draw[BookTeal, fill=BookCodeBg, line width=0.8pt] (10,4.25) rectangle (13,4.75);
  \node at (11.5,4.5) {录 $N{+}3$};
  \node[hw label, anchor=east] at (0.75,3.4) {深度 2 · GPU};
  \draw[BookAccent, fill=BookCodeBg, line width=0.8pt] (7,3.15) rectangle (10,3.65);
  \node at (8.5,3.4) {执行 $N$};
  \draw[BookAccent, fill=BookCodeBg, line width=0.8pt] (10,3.15) rectangle (13,3.65);
  \node at (11.5,3.4) {执行 $N{+}1$};
  \draw[BookMuted, line width=0.6pt, <->] (1,2.75) -- (10,2.75);
  \node[hw label, anchor=north] at (5.5,2.68) {$+2$ 帧 $\approx33$ ms};
  % 深度 3
  \node[hw label, anchor=east] at (0.75,1.9) {深度 3 · CPU};
  \draw[BookTeal, fill=BookCodeBg, line width=0.8pt] (1,1.65) rectangle (4,2.15);
  \node at (2.5,1.9) {录 $N$};
  \draw[BookTeal, fill=BookCodeBg, line width=0.8pt] (4,1.65) rectangle (7,2.15);
  \node at (5.5,1.9) {录 $N{+}1$};
  \draw[BookTeal, fill=BookCodeBg, line width=0.8pt] (7,1.65) rectangle (10,2.15);
  \node at (8.5,1.9) {录 $N{+}2$};
  \draw[BookTeal, fill=BookCodeBg, line width=0.8pt] (10,1.65) rectangle (13,2.15);
  \node at (11.5,1.9) {录 $N{+}3$};
  \node[hw label, anchor=east] at (0.75,0.8) {深度 3 · GPU};
  \draw[BookAccent, fill=BookCodeBg, line width=0.8pt] (10,0.55) rectangle (13,1.05);
  \node at (11.5,0.8) {执行 $N$};
  \draw[BookMuted, line width=0.6pt, <->] (1,0.15) -- (13,0.15);
  \node[hw label, anchor=north] at (7,0.08) {$+3$ 帧 $\approx50$ ms};
  \node[hw label] at (7,-0.7) {深度每加一级：掉帧更平滑，输入时延多一帧——竞技游戏锁深度一，主机游戏常选两到三};
\end{waveform}
\diagramnote{提交深度 1/2/3 的时延—吞吐权衡（60 Hz，横轴每格一帧）。深度 1：CPU 与 GPU 错开一整帧并行，输入到显示约多一帧；深度 2、3：GPU 永远有活干，偶发的慢帧被队列吸收、掉帧更平滑，输入时延相应多两到三帧。深度用内存与时延换吞吐——手快的玩家对 30 ms 以上的差距已有明确体感，这是竞技游戏把深度压到一的全部理由。}

这正引出本节的工程结论：最好的同步，是设计上不需要同步
。帧缓冲、描述表、常量缓冲按帧轮换双份或三份，生产者和
消费者各用各的份，fence 退化为回收水线而非等待点；队列
深度、引擎数量，都是为“互不等待”服务的。反过来，凡是
代码里出现“提交、立即等待、再提交”的段落，基本可以断
定：那里丢掉的吞吐，比省下的内存值钱得多。

真实的等待策略通常是轮询与事件的混合体：先自旋几十微秒
，等到了就赚回中断开销，等不到再登记睡眠——这与第五
章中断合并、存储与外设队列案例轮询模式（如高性能 NVMe 驱动 polling
CQ）是同一种权衡：固定成本与等待时长之比决定谁划算。
功耗维度也值得记一笔：轮询模式把一个核烧在空转上，台式
机上换来最低的时延尚可接受，笔记本与手机上则直接影响续
航；事件模式让出 CPU，核可以降频甚至进深度睡眠。同步模
式的选择从来不只是时延题，也是功耗题。

\topic{显示扫描的固定时序}
渲染引擎按命令执行，显示控制器按显示时钟运行——两者是互相独
立的硬件。显示控制器（scanout engine）不认命令环：它按
像素时钟逐行从帧缓冲读像素，行内先发有效像素、再进入行
消隐；帧内扫完有效行、再进入场消隐，循环往复，与 GPU
忙不忙毫无关系。这套时序是 CRT 时代的遗产：消隐期当年
是电子束回扫的时间，数字接口把整张时序表原样保留下来，
场消隐的空档如今还被用来夹带音频与辅助数据。以
1080p60 的标准时序为例：有效区 $1920\times1080$，含消
隐的总帧 $2200\times1125$，像素时钟 148.5 MHz，每一个
数都能算出物理含义：

\begin{lstlisting}
行时间 = 2200 像素 / 148.5 MHz ≈ 14.8 微秒
帧时间 = 1125 行 x 14.8 微秒 ≈ 16.67 ms（即 1/60 s）
行消隐 = 2200 - 1920 = 280 像素 ≈ 1.9 微秒
场消隐 = 1125 行 - 1080 行 = 45 行 x 14.8 微秒 ≈ 0.67 ms
扫描带宽 = 1920 x 1080 x 4 B x 60 ≈ 0.5 GB/s（4K60 约 2 GB/s）
\end{lstlisting}

三组数字各有工程后果。帧时间 16.67 ms 是渲染的预算线：
一帧的全部工作必须装进这个窗口，装不进就掉帧。场消隐
0.67 ms 是每帧唯一允许“动手脚”的空档：切换帧缓冲、更
新颜色查找表，都要在这个窗口内完成，错过就再等一帧。扫
描带宽 0.5 GB/s 相对 820 GB/s 的显存带宽只是零头，所以
显示扫描从不是带宽问题——它是节拍问题：渲染完成不等于
显示完成，写进帧缓冲的内容要等扫描线走到那一行才真正可
见，从“渲染完”到“用户看见”最多隔将近一帧。显示控制
器的独立性还带来一个副产物：渲染引擎即使完全停止响应，最后
一帧画面照样被一遍遍扫出来，屏幕定格而非黑屏——多屏
机器上每块屏各有一套控制器与像素时钟，各走各的节拍，一
个窗口横跨两屏时，两屏的消隐期也未必对齐。

扫描的顺序本身也值得说细一点：每一帧从左上角逐像素向右
、逐行向下，像读书一样扫完整屏，这个次序叫光栅扫描
（raster scan），图形学里“光栅化”一词正由此来。历史上
还有过隔行扫描的折中：1080i 把一帧拆成奇数行、偶数行两
个场，各扫一次，用一半带宽换到接近逐行的观感——那是
广播带宽紧张年代的产物，今天的接口带宽宽裕，逐行扫描已
经一统天下。接口带宽若真不够，现代标准的答案是显示流压
缩（DSC）：在显示控制器出口处把像素流做视觉无损的轻量
压缩，接口带宽立省三分之二，解压则放在显示器内部——又
是一次“用计算换带宽”，与本章开头 CPU/GPU 面积取舍是同
一条思路。

把行、场两个尺度的时序画在同一张图上，消隐窗口
与交换点的位置关系就清楚了；图中也提前标出下一
题要展开的 tearing。

\begin{waveform}[x=0.8cm,y=0.85cm]
  % 参考线：场消隐边界 / 扫描中途的错误写入点
  \draw[BookRule, line width=0.45pt, dashed] (10.5,4.55) -- (10.5,0.55);
  \draw[BookRule, line width=0.45pt, dashed] (12.5,4.55) -- (12.5,0.55);
  \draw[BookRed, line width=0.45pt, dashed] (5,4.55) -- (5,0.55);
  % 行内时序
  \node[hw label, anchor=east] at (0.5,5.5) {行内};
  \draw[BookTeal, fill=BookCodeBg, line width=0.8pt] (1,5.2) rectangle (8,5.8);
  \node at (4.5,5.5) {有效像素（1920 个）};
  \draw[BookRule, dashed, line width=0.7pt] (8,5.2) rectangle (9.6,5.8);
  \node[hw label] at (8.8,5.5) {行消隐};
  \draw[BookTeal, fill=BookCodeBg, line width=0.8pt] (9.6,5.2) rectangle (12.5,5.8);
  \node at (11.05,5.5) {下一行……};
  \draw[BookMuted, line width=0.6pt, <->] (1,6.12) -- node[hw label, above]{一行 $\approx14.8\ \mu$s} (8,6.12);
  % 帧内时序
  \node[hw label, anchor=east] at (0.5,3.9) {帧内};
  \draw[BookTeal, fill=BookCodeBg, line width=0.8pt] (1,3.6) rectangle (10.5,4.2);
  \node at (7.0,3.9) {有效区：1080 行};
  \draw[BookRule, dashed, line width=0.7pt] (10.5,3.6) rectangle (12.5,4.2);
  \node[hw label] at (11.5,4.48) {场消隐};
  \node[hw label] at (11.5,3.28) {$\approx0.67$ ms};
  % 垂直同步
  \node[hw label, anchor=east] at (0.5,2.6) {垂直同步};
  \draw[BookAccent, line width=0.9pt] (1,2.35) -- (10.5,2.35) -- (10.5,2.85) -- (12.5,2.85);
  % 帧缓冲基址
  \node[hw label, anchor=east] at (0.5,1.3) {帧缓冲基址};
  \draw[BookTeal, fill=BookCodeBg, line width=0.8pt] (1,1.0) rectangle (10.5,1.6);
  \node at (7.2,1.3) {基址 = 前缓冲 A};
  \draw[BookTeal, fill=BookCodeBg, line width=0.8pt] (10.5,1.0) rectangle (12.5,1.6);
  \node[hw label] at (11.5,1.3) {改为 B};
  % 标注
  \draw[BookMuted, line width=0.6pt, <->] (10.5,0.78) -- node[hw label, below]{交换窗口：只在此改写基址} (12.5,0.78);
  \node[hw label, anchor=west] at (5.15,0.15) {红虚线：扫描中途写入 $\Rightarrow$ 上半旧帧、下半新帧（tearing）};
\end{waveform}
\diagramnote{显示扫描的两层时序（示意，不按比例）：行内先发有效像素再入行消隐，一行约 $14.8\ \mu$s；帧内扫完 1080 个有效行进入场消隐，垂直同步脉冲随之拉高。双缓冲的交换——改写帧缓冲基址寄存器——只允许发生在场消隐窗口内；若在扫描中途改写（红虚线），屏幕上半是旧帧、下半是新帧，这就是 tearing。}

\topic{tearing 与双缓冲/垂直同步}
单缓冲把渲染与扫描的矛盾直接摆上台面：渲染引擎写的就是
显示控制器正在读的那块帧缓冲，新内容写进扫描线尚未走到
的区域，这一屏就露出半张新图；画面快速横向运动时，上下
两半错开一道水平缝——这就是 tearing。静态画面的撕裂不
可见，运动越大缝越明显。双缓冲的解法是把读写分开：渲染
写后缓冲，显示读前缓冲，交换只在场消隐那 0.67 ms 内执
行——交换本身只是改写显示控制器的帧缓冲基址寄存器，
一次寄存器写，零拷贝。只要交换点严格避开扫描窗口，撕裂
从机制上消失，这就是垂直同步（vertical sync，vsync）：
把缓冲交换锁在场消隐。

读写分开之后，交换的瞬间动作画出来不过是一次寄存
器写：

\begin{pdfdiagram}
  % 前缓冲：正在被扫描
  \node[hw memory, minimum width=3.4cm, minimum height=1.0cm] (bufA) at (4.3,0) {前缓冲 A};
  \draw[BookRed, line width=0.9pt] (5.6,-0.5) -- (5.6,0.5);
  \node[hw label, anchor=south] at (5.6,0.58) {扫描指针};
  \node[hw compute, minimum width=2.6cm, minimum height=1.0cm] (disp) at (9.3,0) {显示控制器};
  \draw[hw bus] (bufA.east) -- (disp.west);
  \node[hw label, anchor=north] at (7.0,-0.62) {顺序读出送屏幕};
  % 后缓冲：正在被渲染
  \node[hw compute, minimum width=1.7cm, minimum height=0.9cm] (ren) at (1.3,-2.6) {渲染引擎};
  \node[hw memory, minimum width=3.4cm, minimum height=1.0cm] (bufB) at (4.5,-2.6) {后缓冲 B：正在渲染};
  \draw[hw bus] (ren.east) -- (bufB.west);
  % 基址寄存器：交换就是改写它
  \node[hw state, minimum width=2.9cm, minimum height=0.9cm] (base) at (9.3,-2.6) {帧缓冲基址寄存器};
  \draw[hw return] (base.north) -- (disp.south);
  \node[hw label, anchor=east] at (9.15,-1.3) {交换＝改写此寄存器};
  \node[hw label] at (5.5,-3.85) {交换只在场消隐那 0.67 ms 内执行：一次寄存器写，零拷贝，A、B 角色互换};
\end{pdfdiagram}
\diagramnote{双缓冲把渲染与扫描分开：渲染引擎写后缓冲 B，显示控制器经扫描指针读前缓冲 A，两块缓冲互不干扰。交换只是改写显示控制器的帧缓冲基址寄存器——一次寄存器写，零拷贝——且只在场消隐那 0.67 ms 内执行，此后 A、B 角色互换。只要交换点严格避开扫描窗口，撕裂从机制上消失，这就是垂直同步；若扫描中途改写基址，屏幕上半是旧帧、下半是新帧，即 tearing。}

锁住的代价立刻出现：渲染若错过窗口，后缓冲就得再等一整
帧才能上前台。60 Hz 下渲染耗时 17 ms——只比预算多
0.33 ms——帧率不是降到 58，而是直接掉到 30，因为每一
帧都被对齐到第二个消隐窗口；输入时延也随之最多添上
16.7 ms。三缓冲退一步：后缓冲备两份，渲染写完一份不必
等交换就能写另一份，掉帧平滑了，内存与时延各添一份。可
变刷新率（VRR，FreeSync/G-Sync 一类）干脆反转主从：显示
器不再按固定节拍催帧，而是等 GPU 宣布“新帧就绪”才开
始扫下一帧，撕裂与等待同时消除；代价是帧间隔不再均匀，
低帧率时要靠倍帧补偿维持面板刷新，而且整条链路——显卡
、线缆、显示器——都得支持。还要补一句桌面环境的现状：
窗口合成器默认按垂直同步的节拍合成，普通窗口程序其实一
直活在等效 vsync 里，只有全屏独占或显式绕过合成的路径
才有上面这些选择。

\begin{longtable}{L{0.16\linewidth}L{0.12\linewidth}L{0.2\linewidth}L{0.36\linewidth}}
\toprule
方案 & 撕裂 & 输入时延 & 代价与适用 \\
\midrule
无同步 & 有 & 最低 & 基准测试、对画质不敏感的场景 \\\tablerowrule
垂直同步 & 无 & 最多添一帧 & 帧率对齐节拍约数：17 ms 渲染掉到 30 fps \\\tablerowrule
三缓冲 & 无 & 介于两者之间 & 多一份帧缓冲内存，掉帧更平滑 \\\tablerowrule
可变刷新率 & 无 & 接近无同步 & 需显示器支持；帧间隔不均，低帧率靠倍帧 \\
\bottomrule
\end{longtable}

把本节与前一题合看，GPU 系统的两条时钟就齐了：一条是“
能跑多快跑多快”的命令环时钟，吞吐优先；一条是“一微秒
也不商量”的扫描时钟，节拍优先。双缓冲与垂直同步正是两
条时钟之间的翻译器：生产者随意快慢，消费者严格准时，交
换只发生在双方都不看的瞬间。工程上选哪一档同步，本质是
在帧率、时延、画质三者之间挑两个——机制层面，谁也给
不出三全其美。

最后把垂直同步的帧率阶梯画清楚，它是“对齐”二字的直接
后果。交换只能发生在消隐窗口，所以 60 Hz 下实际帧率只能
取节拍的约数：16.7 ms 内完成是 60 fps，33.3 ms 是 30，
50 ms 是 20，66.7 ms 是 15——帧耗时 17 ms 与 33 ms 在
屏幕上得到同一个 30 fps，60 与 30 之间没有任何中间档，
渲染稍慢一档，帧率就跳一档，体感上就是“忽然从 60 掉到
30”。VRR 把阶梯抹平：面板刷新率跟随渲染完成的时刻，17
ms 的帧就以约 59 Hz 显示，23 ms 就以约 43 Hz 显示，帧率
第一次可以连续变化。代价也写在机制里：帧间隔忽长忽短，
对帧间隔敏感的人会觉得“不匀”；低到面板下限（常见 48
Hz）以下时，驱动要靠低帧率补偿把每帧显示两三遍，把名义
刷新率抬回面板能接受的范围。选择依旧回到那句：帧率、时
延、画质，三选二。

把固定刷新与可变刷新画成两条帧时序，阶梯与抹平各是什么
形状就清楚了：

\begin{waveform}[x=0.8cm,y=0.85cm]
  % 组一：固定刷新
  \node[hw label, anchor=west] at (1,7.75) {固定刷新（60 Hz，每格 16.7 ms）};
  \node[hw label, anchor=east] at (0.75,7.0) {vsync 节拍};
  \foreach \x in {1,5,9,13} \draw[BookAccent, line width=0.9pt] (\x,6.75) -- (\x,7.25);
  \node[hw label, anchor=east] at (0.75,5.85) {渲染};
  \draw[BookTeal, fill=BookCodeBg, line width=0.8pt] (1,5.6) rectangle (5.1,6.1);
  \node at (3.05,5.85) {渲染 X：17 ms};
  \draw[BookTeal, fill=BookCodeBg, line width=0.8pt] (9,5.6) rectangle (13.1,6.1);
  \node at (11.05,5.85) {渲染 Y：17 ms};
  \node[hw label, anchor=east] at (0.75,4.7) {上屏};
  \draw[BookAccent, fill=BookCodeBg, line width=0.8pt] (9,4.45) rectangle (13,4.95);
  \node at (11,4.7) {显示 X（等一拍）};
  % 完成迟 0.3 ms：错过当拍窗口，对齐到下一拍
  \draw[BookRed, line width=0.45pt, dashed] (5.1,6.2) -- (5.1,5.5);
  \draw[hw danger] (5.1,5.45) -- (5.1,5.08) -- (8.9,5.08) -- (8.9,5.0);
  \node[hw label, anchor=south, text=BookRed] at (7.5,5.14) {错过窗口：等下一拍（60 掉 30）};
  % 组二：可变刷新
  \node[hw label, anchor=west] at (1,3.8) {可变刷新 VRR：刷新跟着渲染走};
  \node[hw label, anchor=east] at (0.75,3.1) {渲染};
  \draw[BookTeal, fill=BookCodeBg, line width=0.8pt] (1,2.85) rectangle (5.1,3.35);
  \node at (3.05,3.1) {渲染 X：17 ms};
  \draw[BookTeal, fill=BookCodeBg, line width=0.8pt] (5.1,2.85) rectangle (9.2,3.35);
  \node at (7.15,3.1) {渲染 Y：17 ms};
  \draw[BookTeal, fill=BookCodeBg, line width=0.8pt] (9.2,2.85) rectangle (13.3,3.35);
  \node[hw label] at (11.25,3.1) {渲染 Z $\cdots$};
  \node[hw label, anchor=east] at (0.75,1.95) {刷新};
  \foreach \x in {5.1,9.2,13.3} \draw[BookAccent, line width=0.9pt] (\x,1.7) -- (\x,2.2);
  \node[hw label, anchor=east] at (0.75,0.8) {上屏};
  \draw[BookAccent, fill=BookCodeBg, line width=0.8pt] (5.1,0.55) rectangle (9.2,1.05);
  \node at (7.15,0.8) {显示 X};
  \draw[BookAccent, fill=BookCodeBg, line width=0.8pt] (9.2,0.55) rectangle (13.3,1.05);
  \node[hw label] at (11.25,0.8) {显示 Y $\cdots$};
  \node[hw label] at (7,-0.35) {固定刷新：上屏只能取节拍约数（60/30/20）；VRR：帧率连续，代价是帧间隔不匀};
\end{waveform}
\diagramnote{固定刷新与可变刷新的帧时序对照（横轴每格约 16.7 ms）。上方固定刷新：交换只能发生在 vsync 窗口，渲染耗时 17 ms、比预算多 0.3 ms，完成时刻错过当拍窗口，只能对齐到下一拍——帧率从 60 直接掉到 30，中间没有任何档。下方 VRR：显示器不再按固定节拍催帧，等 GPU 宣布帧就绪才扫下一帧，17 ms 的帧按约 59 Hz 显示，帧率第一次可以连续变化；代价是帧间隔忽长忽短，低到面板下限以下时靠倍帧补偿。}

\section{五、从图形管线到通用计算}

本节把前四节装回两种真实负载。

\topic{图形管线的固定与可编程部分}
实时图形几十年只做一件事：把三维场景变成一帧二维像素，每秒至少
60 次。这件事天然分层。顶点处理：给每个顶点的坐标乘上模型、视
图、投影三个矩阵，变换到屏幕空间，顺带算出顶点的颜色与纹理坐标。
图元装配：把变换后的顶点按连接关系拼成三角形，裁掉落在视锥之外
的部分。光栅化：对每个三角形枚举它覆盖哪些像素，在覆盖点上对顶
点属性做插值，输出一批片元（fragment）——片元是“候选像素”，
带着插值出的位置、颜色与纹理坐标。片元着色：对每个片元求最终颜
色，查纹理、算光照。输出合并：把片元与帧缓冲里已有的像素做深度
测试（近的挡住远的）、模板测试与颜色混合，活下来的才写进帧缓冲。
从顶点到像素，数据形态换了四次，每一级的输出就是下一级的输入。

把这五级画成一条流水线，数据形态的接力与可编程
的分界就看得更清楚。

\begin{pdfdiagram}
  \node[hw compute, minimum width=1.65cm, minimum height=1.05cm] (p1) at (0.2,0) {顶点处理\\可编程};
  \node[hw control, minimum width=1.65cm, minimum height=1.05cm] (p2) at (2.3,0) {图元装配\\固定};
  \node[hw control, minimum width=1.65cm, minimum height=1.05cm] (p3) at (4.4,0) {光栅化\\固定};
  \node[hw compute, minimum width=1.65cm, minimum height=1.05cm] (p4) at (6.5,0) {片元着色\\可编程};
  \node[hw control, minimum width=1.65cm, minimum height=1.05cm] (p5) at (8.6,0) {输出合并\\固定};
  \draw[hw bus] (-1.35,0) -- (p1.west);
  \draw[hw bus] (p1.east) -- (p2.west);
  \draw[hw bus] (p2.east) -- (p3.west);
  \draw[hw bus] (p3.east) -- (p4.west);
  \draw[hw bus] (p4.east) -- (p5.west);
  \draw[hw bus] (p5.east) -- (10.35,0);
  \node[hw label] at (-1.0,0.72) {顶点数据};
  \node[hw label] at (1.25,0.72) {变换后顶点};
  \node[hw label] at (3.35,0.72) {三角形};
  \node[hw label] at (5.45,0.72) {片元};
  \node[hw label] at (7.55,0.72) {着色后片元};
  \node[hw label] at (9.9,0.72) {帧缓冲};
  \node[hw label] at (0.2,-0.95) {每顶点};
  \node[hw label] at (2.3,-0.95) {每图元};
  \node[hw label] at (4.4,-0.95) {每三角形};
  \node[hw label] at (6.5,-0.95) {每片元};
  \node[hw label] at (8.6,-0.95) {每像素};
  \node[hw label] at (4.5,-1.6) {蓝框：可编程（shader 在 SM 上按 SIMT 执行）；红框：固定硬件（算法几十年不变，专用电路）};
\end{pdfdiagram}
\diagramnote{图形管线五阶段：数据形态沿顶点、三角形、片元、颜色一路变换，每一级的输出就是下一级的输入。可编程与固定的分界清清楚楚——顶点处理与片元着色交给 shader，在 SM 上按 SIMT 跑；装配、光栅化、输出合并的算法几十年不变，做成固定电路。级间全走先进先出队列，不需要全局调度者。}

哪些可编程，一句话：需要“每个图元各不相同”的数学开放给 shader，
永远同一件事的步骤做成固定硬件。顶点处理与片元着色是可编程的：
游戏要的变换与光照千变万化，厂商把这两级交给程序员写 kernel—
—顶点 shader 与片元 shader，编译成与通用内核一样的指令，在 SM
上按 SIMT 方式跑。图元装配、光栅化、输出合并是固定的：裁剪就是
那几个不等式，光栅化就是逐像素判“点落在三角形哪一侧”的边函数，
深度测试就是一次比较加一次条件写——算法几十年不变，做成专用电
路比可编程核快得多、省得多。裁剪留在固定硬件还有个隐蔽理由：它
的输出不规则，一个三角形裁完可能剩零到几个多边形，不适合 SIMT
的整齐数组风格。

\begin{longtable}{L{0.15\linewidth}L{0.31\linewidth}L{0.20\linewidth}L{0.18\linewidth}}
\toprule
阶段 & 干什么 & 谁来做 & 并行粒度 \\
\midrule
顶点处理 & 矩阵变换、顶点属性计算 & 可编程（顶点 shader） & 每顶点一个线程 \\\tablerowrule
图元装配与裁剪 & 拼三角形、视锥剔除 & 固定硬件 & 每图元 \\\tablerowrule
光栅化 & 枚举覆盖像素、属性插值 & 固定硬件 & 每三角形 \\\tablerowrule
片元着色 & 纹理采样、光照、最终颜色 & 可编程（片元 shader） & 每片元一个线程 \\\tablerowrule
输出合并 & 深度与模板测试、混合写回 & 固定硬件 & 每像素 \\
\bottomrule
\end{longtable}

再算一笔负载账，看面积该怎么分。4K 帧缓冲
$3840\times2160\approx8.3\times10^{6}$ 像素，60 Hz 下每秒产出
约 $5\times10^{8}$ 像素；画面里三角形层层遮挡，过绘制按保守 3
倍计，片元着色每秒要处理约 $1.5\times10^{9}$ 个片元。每片元做
纹理过滤加几项光照按 100 次浮点操作估，就是每秒
$1.5\times10^{11}$ 次操作，还没算顶点与光栅化；CPU 单核每秒能
做的浮点操作在 $10^{10}$ 到 $10^{11}$ 量级（标量到向量峰值），连片元着色这一级也无法满足吞吐要求。这就
是图形管线必须并行、且把并行做成流式的理由：片元数量比顶点大一
到两个数量级，面积自然跟着数据量走——可编程的片元级占去芯片大
部分，固定级做窄而深。

把片元着色里最重的一次访存——纹理采样——单独画出来，纹
理为什么配专用 cache 与二维布局就有了形状：

\begin{pdfdiagram}
  \node[hw memory, minimum width=2.2cm] (vram) at (1.1,0) {VRAM 纹理区\\二维分块布局};
  \node[hw memory, minimum width=1.7cm] (tc) at (3.9,0) {纹理\\cache};
  \node[hw compute, minimum width=2.1cm] (tu) at (6.6,0) {纹理单元\\寻址＋插值过滤};
  \node[hw compute, minimum width=2.3cm] (fs) at (9.45,0) {片元着色器\\（SM，SIMT）};
  \draw[hw bus] (vram.east) -- (tc.west);
  \draw[hw bus] (tc.east) -- (tu.west);
  \draw[hw bus] (tu.east) -- (fs.west);
  \draw[hw bus] (fs.east) -- (11.35,0);
  \node[hw label, anchor=south west] at (10.75,0.12) {颜色};
  % 未命中：按二维邻近取回整行
  \draw[hw return] (tc.north) -- (3.9,1.0) -- (1.1,1.0) -- (vram.north);
  \node[hw label, anchor=south] at (2.5,1.05) {未命中：按二维邻近取回整行};
  \node[hw label, anchor=north] at (2.6,-0.5) {命中：纹素};
  \node[hw label, anchor=north] at (5.15,-0.5) {双线性过滤};
  \node[hw label, anchor=north] at (8.0,-0.5) {纹素＋光照};
  \node[hw label] at (5.5,-1.6) {同一小片纹理一帧内被成千上万个像素反复采样，命中频次以亿计——读多写少};
  \node[hw label] at (5.5,-2.2) {相邻像素采相邻纹素：取回的整行下一拍还能用（二维局部性）};
\end{pdfdiagram}
\diagramnote{一次纹理采样的路径。片元着色器发起采样，纹理单元算地址、做插值过滤；命中纹理 cache 直接得纹素，未命中才按二维邻近性从 VRAM 取回整行——纹理在显存里按二维分块布局存放，相邻像素采相邻纹素，取回的整行下一拍还能用。同一小片纹理一帧内被成千上万个像素反复采样，命中频次以亿计：读多写少加二维局部性，是纹理配专用 cache 的全部理由。}

更妙的是这条管线不需要调度者。每级只盯着自己的输入队列：顶点处
理完一个顶点就交给装配，光栅化每产出一个片元就传递给片元着色器；
级间是先进先出队列，级内是成千上万个互不相同的顶点与片元。数据
并行叠着流水并行，没有全局汇合，没有全局同步，每一级按自己的节
奏铺满硬件——流式并行是图形的天然结构，也是后来通用计算直接继
承的结构。

\topic{一次绘制命令的旅行}
把“应用画一个三角形”从 API 调用到屏幕发光按时序走一遍。这是
第四节的命令环与显示扫描在图形负载上的一次完整应用，这里只追
数据流，同步原语的细节不再重复。

\tablechunkintro{1}{2}{1}{7}
\begin{longtable}{L{0.05\linewidth}L{0.63\linewidth}L{0.16\linewidth}}
\toprule
步 & 动作 & 执行者 \\
\midrule
1 & 应用把顶点数据（坐标、颜色、纹理坐标）放入显存可见的缓冲，调用绘制命令 & CPU \\\tablerowrule
2 & 驱动校验渲染状态，把绘制翻译成 GPU 命令包，写入命令环尾部 & CPU \\\tablerowrule
3 & 屏障之后写 doorbell 寄存器（MMIO 写），通知 GPU 有新命令 & CPU \\\tablerowrule
4 & 命令处理器 DMA 取走命令包，解析出绘制，配置管线状态 & GPU 前端 \\\tablerowrule
5 & 顶点 fetch：DMA 按索引从显存读出全部顶点的属性 & GPU 前端 \\\tablerowrule
6 & 顶点着色：每个顶点一个线程，在 SM 上跑顶点 shader & SM \\\tablerowrule
7 & 图元装配与裁剪：拼出三角形，确认落在视锥之内 & 固定硬件 \\
\bottomrule
\end{longtable}

\tablechunkintro{2}{2}{8}{12}
\begin{longtable}{L{0.05\linewidth}L{0.63\linewidth}L{0.16\linewidth}}
\toprule
步 & 动作 & 执行者 \\
\midrule
8 & 光栅化：枚举三角形覆盖的像素，插值属性，产出一批片元 & 固定硬件 \\\tablerowrule
9 & 片元着色：每个片元一个线程，查纹理、算光照，得到颜色 & SM \\\tablerowrule
10 & 输出合并：深度测试、混合，写帧缓冲 & 固定硬件 \\\tablerowrule
11 & 显示控制器按行扫描时序从帧缓冲读像素，逐行送出 & 显示引擎 \\\tablerowrule
12 & 像素经显示接口抵达屏幕，对应位置发光 & 显示器 \\
\bottomrule
\end{longtable}

这张表有三处值得停下来看。第一，第 1 到第 3 步与存储与外设队列案例中的 NVMe
的提交一模一样：备好数据、写命令、doorbell——存储与设备事务相关章节“描述符先
可见、再通知”的约定原样适用。第二，CPU 在第 3 步就结束了本次
绘制的参与：它不等待渲染完成，转身去准备下一帧的命令；GPU 从
第 4 步起自主跑完整个管线。命令环换来的正是这份解耦：CPU 与
GPU 各跑各的，只在环的生产与消费处碰头。第三，第 11 步与前面
所有步骤异步：显示扫描按自己的固定时钟走，读的是“此刻完整的
帧缓冲”，渲染完成与显示完成之间隔着第四节讲的垂直同步与帧切换。
一帧 16.7 ms 的预算里，CPU 备命令通常只占一二毫秒，其余都是
GPU 的管线时间；哪一级超载哪一级的队列先满，帧率被最慢的一级
锁死——管线吞吐由瓶颈级决定，这与“冒险、异常与性能”一章的利用率分析是同一条
规律。

\topic{CUDA/OpenCL 的内核模型}
通用计算借用同一台机器，只换一层软件抽象。内核（kernel）就是
“在 N 个数据点上运行的同一个函数”：程序员写处理一个元素的代
码，启动时给出总数，硬件负责把 N 个实例铺开。CUDA 把一次启动
组织成三层：网格（grid）是这次启动的全部线程；块（block）是网
格切成的组，块内线程可以经共享内存协作、可以块内同步；线程是
最小单位，拿着自己的索引算自己那份。三层到硬件的映射很直接：
块整体调度到一个 SM 上常驻；块内每 32 个线程编成一条 warp；
warp 是 SM 真正发射指令的单位，一条 lane 对应一个线程。OpenCL
的 NDRange、work-group、work-item 是同一张表的另一套名词。向
量加法内核的骨架如下：

\begin{lstlisting}
/* 向量加法：c[i] = a[i] + b[i]，共 N 个元素 */
__global__ void vec_add(const float *a, const float *b,
                        float *c, int N) {
    int i = blockIdx.x * blockDim.x + threadIdx.x;  /* 全局索引 */
    if (i < N)            /* 越界的线程直接退出 */
        c[i] = a[i] + b[i];
}

/* 主机侧：N = 2^20，每块 256 线程 */
int threads = 256;
int blocks  = (N + threads - 1) / threads;   /* 4096 块 */
vec_add<<<blocks, threads>>>(d_a, d_b, d_c, N);
\end{lstlisting}

算一下这次启动的规模。$N=2^{20}=1048576$，每块 256 线程，得
4096 块；每块 $256/32=8$ 条 warp。若 GPU 有 16 个 SM、每 SM 可
常驻 8 块，则在途 128 块、$128\times256=32768$ 个线程——本章
开头说的“上万线程”就是这样来的：在途的只是第一梯队，约占总线
程数的 3\%，其余近四千块在队列里候场，某块跑完立即补位。内核
里“越界即退出”那行判断，是网格与 N 不整除时的边界保护（本例
恰好整除，但写成习惯）。再注意内核代码里没有循环、没有同步、
没有“第几个元素归谁”的概念：并行度全部交给启动配置。这与第四
章最小 CPU“一条指令流顺序走”是两种世界观：程序员交出了控制流，
换回来的是上万份数据同时推进。

三层软件名词到三层硬件的对应画成对照，这次启动的
规模也就摆齐了：

\begin{pdfdiagram}
  \node[hw state, minimum width=4.2cm] (kg) at (2.3,0) {网格 grid / NDRange\\4,096 块};
  \node[hw compute, minimum width=5.6cm] (ksm) at (8.0,0) {16 个 SM：块整体入驻、常驻\\在途 $16\times8=128$ 块};
  \draw[hw bus] (kg.east) -- (ksm.west);
  \node[hw label, anchor=south] at (4.8,0.52) {调度};
  \node[hw state, minimum width=4.2cm] (kb) at (2.3,-1.25) {块 block / work-group\\256 线程 $=$ 8 条 warp};
  \node[hw compute, minimum width=5.6cm] (kw) at (8.0,-1.25) {warp：SM 发射指令的单位\\块内每 32 个线程编一组};
  \draw[hw bus] (kb.east) -- (kw.west);
  \node[hw label, anchor=south] at (4.8,-0.73) {编组};
  \node[hw state, minimum width=4.2cm] (kt) at (2.3,-2.5) {线程 thread / work-item\\拿自己的索引算自己那份};
  \node[hw compute, minimum width=5.6cm] (kl) at (8.0,-2.5) {lane：一条算术道\\对应组内一个线程};
  \draw[hw bus] (kt.east) -- (kl.west);
  \node[hw label, anchor=south] at (4.8,-1.98) {绑定};
  \node[hw label] at (5.5,-3.4) {在途 128 块 $=32{,}768$ 线程，约占总线程数的 3\%；其余近四千块在队列里候场，跑完一块补一块};
\end{pdfdiagram}
\diagramnote{CUDA 三层抽象到硬件的映射（OpenCL 名词列在斜线后）。网格是这次启动的全部线程，切成 4,096 块；块整体调度到一个 SM 上常驻，16 个 SM 各驻 8 块，在途 128 块合 32,768 个线程；块内每 32 个线程编成一条 warp，warp 才是 SM 真正发射指令的单位，一条 lane 对应一个线程。程序员只写处理一个元素的代码，并行度全部交给启动配置。}

\topic{案例：一次内核启动的完整 trace}
把上面这个 \code{vec\_add} 从数据到结果走完。设 $N=2^{20}$，
GPU 接在 PCIe 3.0 x16 上（单向约 15.75 GB/s），显存带宽按 400
GB/s 计，输入 a、b 共 8 MiB，输出 c 共 4 MiB。

\begin{longtable}{L{0.05\linewidth}L{0.50\linewidth}L{0.14\linewidth}L{0.15\linewidth}}
\toprule
步 & 动作 & 耗时量级 & 执行者 \\
\midrule
1 & 把 a、b 放入锁页（pinned）缓冲（存储与设备事务相关章节：DMA 要求物理地址稳定） & 一次性 & CPU \\\tablerowrule
2 & 发拷贝命令，DMA 把 8 MiB 经 PCIe 搬进显存 & 约 530 $\mu$s & GPU DMA \\\tablerowrule
3 & 启动命令（网格配置、参数指针）写入命令环，写 doorbell & 微秒级 & CPU \\\tablerowrule
4 & GPU 取命令，把 4096 块分发到各 SM & 微秒级 & GPU 前端 \\\tablerowrule
5 & 内核执行：读 a、b 写 c，显存流量共 12 MiB & 约 31 $\mu$s & SM \\\tablerowrule
6 & 全部块完成，GPU 把 fence 写到主存并发中断 & 微秒级 & GPU \\\tablerowrule
7 & DMA 把 4 MiB 结果搬回主存 & 约 266 $\mu$s & GPU DMA \\\tablerowrule
8 & fence 就绪，CPU 读结果缓冲 & 微秒级 & CPU \\
\bottomrule
\end{longtable}

总时长约 830 $\mu$s。拆开看占比：两次 PCIe 拷贝约 800 $\mu$s，
占九成半；内核本身 31 $\mu$s，只占 4\%。根源是算术强度低得无
可救药：全内核只做 $2^{20}$ 次加法，流量 12 MiB，约 0.08 次操
作每字节；按“冒险、异常与性能”一章 roofline，峰值算力 10 TFLOPS、带宽 400 GB/s
的机器，拐点在 $10^{13}/(4\times10^{11})=25$ 次操作每字节，这
个内核差了近三百倍，注定受带宽上限约束。工程结论有两条。其一，
向量加法这种强度的负载不值得上 GPU：数据还没算完，搬运先花掉
二十几倍的时间；值得上 GPU 的负载，必须让每个字节多做几十上百
次操作。其二，真要用 GPU 就得批：数据常驻显存，一次搬运喂一串
内核，把 PCIe 这笔固定开销摊到足够多的计算上——分块流水、异步
重叠，都是在跟这九成半作战。

把总时长按真实比例画成一条时间条，钱花在哪一头清晰可见：

\begin{pdfdiagram}
  % 三段按真实比例：530 + 31 + 266 ≈ 830 µs
  \filldraw[fill=BookAccent!18, draw=BookAccent, line width=0.5pt] (0.4,0) rectangle (7.05,0.7);
  \node[hw label, text=BookInk] at (3.7,0.35) {拷入 8 MiB $\approx$530 $\mu$s};
  \filldraw[fill=BookTeal!30, draw=BookTeal, line width=0.5pt] (7.05,0) rectangle (7.44,0.7);
  \draw[BookMuted, line width=0.5pt] (7.25,0.7) -- (7.25,1.25);
  \node[hw label, anchor=south] at (7.25,1.3) {内核 31 $\mu$s（约 4\%）};
  \filldraw[fill=BookAmber!20, draw=BookAmber, line width=0.5pt] (7.44,0) rectangle (10.79,0.7);
  \node[hw label, text=BookInk] at (9.1,0.35) {拷回 4 MiB $\approx$266 $\mu$s};
  \draw[BookMuted, line width=0.6pt, <->] (0.4,-0.3) -- (10.79,-0.3);
  \node[hw label, anchor=north] at (5.6,-0.4) {总时长 $\approx$830 $\mu$s};
  \node[hw label] at (5.6,-1.15) {PCIe 搬运两段合计 $\approx$796 $\mu$s，占九成半——优化指向批量化与数据常驻，不在内核本身};
\end{pdfdiagram}
\diagramnote{一次内核启动的耗时构成（$N=2^{20}$ 向量加，PCIe 3.0 x16，按比例）：拷入 8 MiB 约 530 $\mu$s，内核执行约 31 $\mu$s，拷回 4 MiB 约 266 $\mu$s，总时长约 830 $\mu$s。两次 PCIe 搬运占九成半，内核只占约 4\%——算术强度约 0.08 FLOP/B 的负载，时间几乎全花在搬运上；工程上要么批量化、让数据常驻显存把这笔固定开销摊薄，要么干脆留在 CPU 上。}

把这次 trace 与存储与外设队列案例中的 NVMe 读并排摆开，是同一个骨架穿了两
套衣服：

\begin{longtable}{L{0.12\linewidth}L{0.38\linewidth}L{0.34\linewidth}}
\toprule
维度 & 内核启动 & NVMe 读（存储与外设队列案例） \\
\midrule
提交物 & 命令包：网格配置、参数指针 & 64 B 命令项：LBA 与 PRP 列表 \\\tablerowrule
队列 & 命令环（第四节） & SQ/CQ 队列对 \\\tablerowrule
通知 & doorbell MMIO 写 & doorbell MMIO 写 \\\tablerowrule
搬运 & DMA 双向：输入进显存、结果回主存 & DMA 单向：NAND 数据进页框 \\\tablerowrule
完成 & fence 写主存加中断 & CQ 完成项加 MSI \\\tablerowrule
大头 & PCIe 搬运，占九成半 & 介质读，占八成 \\
\bottomrule
\end{longtable}

骨架相同不是巧合：凡是“CPU 提交命令、设备执行”的接口，都要回答同
样四个问题——命令放哪（内存队列）、怎么通知（doorbell）、数
据谁搬（DMA）、完成怎么确认（完成项或 fence）。NVMe 与 GPU 驱
动是这套约定的两个实例，学会一个，另一个的文档就只剩名词表要
背。差别在优化的主战场：NVMe 的大头是介质，优化指向访问形态
（对齐、合并、足够的在途深度）；GPU 内核的大头是搬运，优化指
向批量化与数据常驻。

\topic{GPU 编程的三个常见陷阱}
GPU 的宽容度很高：内核写糙了照样算出正确结果，性能却能差出百
倍。最常见的三个陷阱看着互不相关，破坏的却是同一个根本——海
量在途工作。

陷阱一：分支分歧。一条 warp 共享一条指令流，遇到各 lane 去向
不同的分支，硬件只能串行：先放行走向 if 一侧的 lane、其余停等，
再放行走向 else 一侧的 lane、其余停等。最坏情形是每条 lane 走
不同的路：

\begin{lstlisting}
/* 反例：按 lane 号散到 32 个分支，warp 退化成顺序执行 */
switch (threadIdx.x % 32) {
    case 0:  c[i] = f0(a[i]);  break;
    case 1:  c[i] = f1(a[i]);  break;
    /* ... 直到 case 31 */
}
\end{lstlisting}

后果直接按 SIMT 模型算：32 条路径串行，这条 warp 的吞吐退到峰
值的 1/32，等于把 32 条 lane 的硬件用成 1 条。随机数据上的普通
if-else 也好不到哪去：约半数 lane 走 A、半数走 B，两条路径都
得执行，耗时是 $t_A+t_B$ 的串行和，而不是两者之一。

陷阱二：未合并访存。一条 warp 的 32 次访存若地址连续，恰好拼
成 128 字节，一次内存事务取回一整条缓存线，人人有份；若地址跨
步跳跃，每个线程各碰一条线：

\begin{lstlisting}
/* 反例：跨步 32，每条 lane 各要一条 128 B 的线 */
c[i] = a[i * 32] + b[i * 32];
\end{lstlisting}

原来一次事务办完的事变成 32 次，每次取 128 字节只用 4 字节，
带宽效率 1/32：400 GB/s 的显存被用成 12.5 GB/s。内核从带宽受
限变成事务受限，内存队列里堆满等数据的 warp——占用率再高，也
藏不住被放大了 32 倍的时延。

陷阱三：CPU-GPU 过度同步。每次启动后立刻等 fence、立刻把结果
拷回、立刻在 CPU 上检查再决定下一步：

\begin{lstlisting}
/* 反例：1000 次小启动，每次都同步等结果 */
for (int k = 0; k < 1000; ++k) {
    vec_add<<<1, 256>>>(d_a + k*256, d_b + k*256,
                        d_c + k*256, 256);
    gpu_synchronize();          /* 等 fence，GPU 排空 */
    check_on_cpu(d_c + k*256);  /* 立刻读回 */
}
\end{lstlisting}

每次同步都把整条管线抽干：doorbell、执行、fence、中断，一趟往
返按 10 $\mu$s 估，1000 次就是 10 ms；同样的数据量一次启动只
需几十微秒，慢出去两个数量级。更要命的是同步期间 GPU 完全空闲：
在途工作清零，时延没有任何东西可藏。三个陷阱的药方其实是同一
句话：让分支收敛回整块整块地走，让访存回到 128 字节一条线的连续模式，
让同步退到一批工作的末尾——都是在把“在途”二字养回来。

\topic{GPU 之外的加速器}
GPU 相对 CPU 已经专了一级：放弃单线程时延，换吞吐。这条路再往
前走是更专的芯片，规律始终同一条：通用性再降一档，效率再升一
档。两个代表。其一是 NPU/TPU 一类的矩阵引擎。神经网络的核心操
作是大矩阵乘加，TPU v1 干脆把 $256\times256$ 个乘加器排成脉动
阵列（systolic array）：数据从阵列一端泵入，在相邻单元间逐级
传递，每拍完成 65536 次乘加，700 MHz 下约 $4.6\times10^{13}$
次乘加每秒（约 92 TOPS）。关键不在运算更快，而在搬运更少：数
据在阵列内部的寄存器间流动，几百次乘加才需要一次对外访存，而
访存正是能耗大头。其二是视频编解码器。H.264/HEVC 的帧间预测、
变换、熵编码是一二十年基本不变的固定流水线，做成专用电路后，手机
SoC 上一小块逻辑就能实时解 4K60 的视频，功耗以瓦计；同样的码
流交给 CPU 软件解码，要动用多个大核，功耗高出一个数量级。

\begin{longtable}{L{0.16\linewidth}L{0.28\linewidth}L{0.24\linewidth}L{0.18\linewidth}}
\toprule
芯片 & 固定下来的计算 & 放弃的通用性 & 换来的效率 \\
\midrule
GPU & SIMT 数组并行、图形管线 & 单线程时延、乱序投机 & 每瓦吞吐高一个量级 \\\tablerowrule
NPU/TPU & 矩阵乘加脉动阵列 & 只剩张量算子可表达 & 每瓦算力再升一档 \\\tablerowrule
视频编解码器 & 预测、变换、熵编码整条流水线 & 几乎不可编程 & 瓦级实时 4K \\
\bottomrule
\end{longtable}

这张表从上往下看，就是加速器的一般规律：每固定一层计算，可编
程性退一档，单位能耗能做的事上一个数量级。CPU 在表外最上头：
什么都能干，什么都最贵。系统设计的问题从来不是“用不用加速器”，
而是“负载里哪一段足够固定，固定到值得为它造一块电路”。这正是
“冒险、异常与性能”一章定量方法给出的决策顺序：先测负载各段占多少时间，再决定
把哪一段交给专用硬件。

把这张表摊成坐标，“通用性换效率”就是一条一眼可见
的走向：

\begin{pdfdiagram}
  % 坐标轴
  \draw[BookMuted, line width=0.6pt, ->] (0.9,0.5) -- (10.5,0.5);
  \node[hw label, anchor=north] at (10.4,0.38) {通用性};
  \draw[BookMuted, line width=0.6pt, ->] (0.9,0.5) -- (0.9,4.9);
  \node[hw label, anchor=south] at (0.9,5.0) {每瓦效率};
  % 四类芯片
  \node[hw state, minimum width=2.4cm, minimum height=0.9cm] (acpu) at (9.2,1.35) {CPU\\什么都能干};
  \node[hw compute, minimum width=2.4cm, minimum height=0.9cm] (agpu) at (6.6,2.3) {GPU\\SIMT 数组并行};
  \node[hw compute, minimum width=2.4cm, minimum height=0.9cm] (anpu) at (4.0,3.25) {NPU/TPU\\矩阵乘加脉动阵列};
  \node[hw memory, minimum width=2.4cm, minimum height=0.9cm] (avc) at (2.35,4.2) {编解码器\\固定流水线};
  \node[hw label] at (7.0,4.6) {每固定一层计算：可编程性退一档，效率升一档};
\end{pdfdiagram}
\diagramnote{通用性—效率坐标上的四类芯片。CPU 在最右下：什么都能干，什么都最贵；GPU 固定了 SIMT 数组并行与图形管线，放弃单线程时延，每瓦吞吐高一个量级；NPU/TPU 只剩张量算子可表达，数据在脉动阵列内部流动、几百次乘加才一次对外访存，每瓦算力再升一档；视频编解码器把整条预测、变换、熵编码流水线做成电路，几乎不可编程，换来瓦级实时 4K。系统设计的问题从来不是“用不用加速器”，而是负载里哪一段固定到值得为它造一块电路。}




























\section{六、CPU 与 GPU 共享地址不等于共享可见性}

集成 GPU 常与 CPU 共享物理内存控制器，独立显卡则拥有本地 VRAM，并通过 PCIe 与主存交换数据。两种平台都可能向程序提供统一虚拟地址，但“同一个指针可被两侧解释”只解决地址命名，不能自动解决页驻留、Cache 一致性、访问权限和完成顺序。运行时必须知道当前页面位于主存还是 VRAM、哪一侧持有可写副本，以及在所有权转换前有哪些写仍停留在 Cache、写合并缓冲或 DMA 队列中。

CPU 准备输入缓冲时，先写普通 Cache line，再由驱动或运行时执行平台要求的同步操作，使数据对设备可见；随后提交命令描述符，执行发布屏障，最后写 doorbell。GPU 只有在命令处理器取到稳定描述符和资源表后才能执行。内核完成计算并写结果后，fence 值的更新是 GPU 向主机发布完成的边界；CPU 等待 fence 成功后还要按平台一致性规则取得最新数据，才能读取结果。

若系统提供硬件一致性，探测协议可以使 CPU 与 GPU Cache 对同一 Cache line 保持一致，但 fence 仍然必要：一致性回答单个内存位置最终看到哪个值，fence 回答一批命令和结果何时全部完成。若系统不提供一致性，运行时需要显式 flush、invalidate 或使用不可缓存/一致性映射。把“总线支持一致性”理解为“不再需要同步”，会让 CPU 在 GPU 尚未结束时读取一个完全一致但尚未更新的旧值。

\begin{longtable}{L{0.18\linewidth}L{0.28\linewidth}L{0.34\linewidth}L{0.12\linewidth}}
\toprule
边界 & 必须稳定的状态 & 可见结果 & 典型错误 \\
\midrule
CPU 发布输入 & 数据、描述符、资源地址与权限 & GPU 读取同一版本输入 & doorbell 早于描述符可见 \\\tablerowrule
地址转换 & 页表、PASID/进程身份、IOMMU/设备 TLB & 设备地址命中允许的物理页 & 页面已回收但设备仍缓存翻译 \\\tablerowrule
GPU 执行 & 命令状态、线程块、VRAM/共享内存 & 结果写入目标缓冲 & 内核越界或设备异常 \\\tablerowrule
GPU 发布完成 & 结果写入顺序与 fence 值 & CPU/后续队列可确认完成 & fence 更新早于结果可见 \\\tablerowrule
显示扫描 & 帧缓冲地址、格式、stride 与翻页时刻 & 显示控制器扫描完整帧 & 渲染与扫描使用不同帧代际 \\\tablerowrule
复位恢复 & 队列 generation、映射与在途 DMA & 旧完成不能污染新任务 & 复位后沿用旧指针和 fence \\
\bottomrule
\end{longtable}

\section{七、页面迁移、设备缺页与超额使用}

统一内存可以在首次访问时把页面迁到使用者附近。GPU 访问尚未驻留的页时，设备 MMU 产生 fault，暂停相关 wave/warp，把地址和访问类型报告给主机；操作系统或驱动分配目标显存、复制页面、更新页表并失效设备 TLB，随后重放访问。这个过程把普通 load 的延迟从纳秒扩展到微秒甚至更高，因此内核即使计算量不大，也可能因频繁迁移而出现明显尾延迟。

若活跃工作集超过 VRAM，页面会在主存与显存间往返。迁移粒度通常是页，而线程访问粒度可能只是几十字节；离散访问会放大 PCIe 流量。运行时可以预取页面、固定关键缓冲、批量安排相邻工作或让部分数据直接驻留主存，但每种策略都要量化带宽与容量。所谓“统一”减少了显式复制代码，不会消除物理链路和容量差异。

设备缺页还必须与 CPU 页框生命周期协调。CPU 不能在 GPU 仍持有翻译或在途访问时回收页面；驱动需要先停止新提交，等待 fence，失效设备页表/IOMMU 映射并确认失效完成，再释放页框。只更新 CPU 页表而遗漏设备 TLB，会让 GPU 使用旧物理地址写入已经分配给其他进程的页面。

\section{八、命令、计算、显示是三个完成域}

图形程序经常把“命令已提交”“渲染已完成”和“画面已显示”混为一件事。CPU 把命令写入环并敲 doorbell，只表示 GPU 可以开始读取；GPU fence 完成表示对应写入达到设备定义的可见边界；显示控制器在垂直消隐期采用新帧缓冲并扫描最后一个像素，才表示该帧真正出现在接口输出上。三者使用不同队列和时钟域。

双缓冲用 front buffer 供显示扫描、back buffer 供 GPU 渲染。GPU 完成 back buffer 后，显示引擎只能在规定翻页边界交换角色，否则屏幕上半部可能来自旧帧、下半部来自新帧。三缓冲允许 GPU 在显示等待垂直同步时继续渲染，但增加一帧排队可能提高输入到显示的时延。选择缓冲数量时应同时看吞吐、撕裂和交互延迟，不能只看平均帧率。

故障定位也应按三个域读取证据：命令读取指针不动，先查队列、doorbell 和设备状态；计算 fence 不前进，查内核异常、资源映射和 GPU fault；fence 已完成但画面不更新，查翻页队列、帧缓冲格式、显示时序与链路。沿这个顺序可以避免把显示扫描问题误判为着色器计算错误。




























\section{九、显示物理通路：从帧缓冲到屏幕发光}

GPU 把像素写入帧缓冲，只完成了“画面已经生成”。显示控制器还要在固定刷新节拍下读取像素，完成平面合成、色彩转换和时序生成；链路控制器再把像素变成 HDMI、DisplayPort 或 eDP 上的高速串行数据；显示器内部的时序控制器最后驱动面板行列电路。直到对应像素真正改变透光率或发光强度，一帧图像才走完物理通路。

\begin{pdfdiagram}[x=1cm,y=1cm]
  \node[hw memory, minimum width=2.5cm] (fb) at (1.3,3.7) {帧缓冲\\像素、stride、格式};
  \node[hw control, minimum width=2.7cm] (dc) at (4.4,3.7) {显示控制器\\取数、合成、时序};
  \node[hw compute, minimum width=2.7cm] (tx) at (7.7,3.7) {链路发送端\\编码、分包、SerDes};
  \node[hw block, minimum width=2.7cm] (link) at (11.0,3.7) {HDMI/DP/eDP\\差分通道与辅助通道};
  \draw[hw bus] (fb) -- (dc);
  \draw[hw bus] (dc) -- (tx);
  \draw[hw bus] (tx) -- (link);

  \node[hw control, minimum width=2.7cm] (rx) at (11.0,1.0) {接收端与 TCON\\恢复时钟、重排像素};
  \node[hw block, minimum width=2.7cm] (panel) at (7.7,1.0) {面板行列驱动\\栅极/源极或像素电流};
  \node[hw state, minimum width=2.7cm] (light) at (4.4,1.0) {LCD/OLED 像素\\透光或主动发光};
  \node[hw state, minimum width=2.5cm] (seen) at (1.3,1.0) {可见画面\\亮度、色彩、刷新};
  \draw[hw bus] (link) -- (rx);
  \draw[hw bus] (rx) -- (panel);
  \draw[hw bus] (panel) -- (light);
  \draw[hw bus] (light) -- (seen);
\end{pdfdiagram}
\diagramnote{一帧画面的完整物理通路。上排从帧缓冲读取像素并转换成串行链路数据，下排从接收、时序控制和行列驱动一直走到像素发光。渲染 fence、翻页、链路传输、面板扫描和人眼看到画面是五个不同的完成边界。}

\topic{显示控制器按扫描节拍消费帧缓冲}

帧缓冲不仅是一片像素数组，还带有宽高、像素格式、每行跨度、平面地址、色域和传递函数。显示控制器通常拥有多个硬件平面：主画面、光标和视频层可以分别取数，再由合成器缩放、混合并转换到链路需要的色彩格式。若一行像素需要的内存带宽在截止时刻前没有得到满足，扫描逻辑不能停下来等待，它会重复旧数据、显示底色或报告 underflow；因此显示带宽必须按像素时钟、每像素位数、平面数量和缩放开销预留，而不能只看 GPU 的平均显存带宽。

刷新时序包含有效像素区与水平、垂直消隐区。现代数字链路不再需要电子束回扫，但发送端与接收端仍用这套时序约定一帧和一行的边界。系统在垂直消隐附近更新 plane 地址，可以让下一帧从新缓冲区开始，避免同一刷新周期混入两代画面。可变刷新率允许延长帧间间隔，却仍要求每一行一旦开始就按约定速率送完。

\topic{HDMI、DisplayPort 与 eDP 先建立能力和链路，再发送视频}

HDMI 通过 DDC 读取接收端的 EDID，了解分辨率、刷新率、色彩格式和音频能力；随后发送端选择双方都支持的时序与链路速率。DisplayPort 使用 AUX 辅助通道读取能力并执行链路训练：接收端反馈时钟恢复和均衡状态，发送端逐步调整电压摆幅与预加重，直到各 lane 在目标速率下可靠工作。eDP 把相同的分组化链路用于设备内部面板，并增加面板电源、背光和低功耗状态的协同。

链路带宽预算必须基于实际传输格式。分辨率和刷新率给出每秒有效像素数，消隐、编码和分包带来额外开销；RGB 8 bit 每分量与 10 bit 每分量、全色度与色度抽样也会改变数据率。若模式超出通道数量与每 lane 速率的组合，系统只能降低刷新率、降低色深、使用压缩或拒绝该模式。所谓“接口支持 4K”没有完整含义，必须同时给出刷新率、色深、色度格式和链路配置。

\topic{TCON 与面板把数字像素落实为电压、电流和光}

接收端恢复串行时钟和数据，重建像素流并交给 TCON（timing controller，时序控制器）。TCON 生成逐行、逐列的选通信号。LCD 像素本身不发光：背光先产生光，薄膜晶体管给液晶单元充电，液晶改变偏振状态，从而控制穿过彩色滤光片的光量；电容保持像素电压直到下次刷新。OLED 子像素则由驱动晶体管控制电流并主动发光，不需要统一背光，但亮度均匀性、老化补偿和低灰阶控制更复杂。

面板响应不是瞬时的。LCD 液晶转动、OLED 驱动与补偿、背光调制和图像处理都可能增加延迟。高动态范围还要求内容元数据、色调映射、面板峰值亮度和局部调光共同配合。链路已经无误并不保证画面正确：偏色可能来自色彩矩阵或传递函数，闪屏可能来自链路训练或面板供电，横线可能来自 TCON、行驱动或连接排线，必须沿边界逐级定位。

\begin{longtable}{L{0.18\linewidth}L{0.34\linewidth}L{0.36\linewidth}}
\toprule
可观察事件 & 它真正证明的事实 & 仍未证明的后续边界 \\
\midrule
GPU fence 完成 & 渲染写入达到设备规定的可见点 & 新帧已被显示控制器采用 \\\tablerowrule
翻页事件完成 & plane 地址已在约定刷新边界切换 & 链路与面板已显示最后一个像素 \\\tablerowrule
链路训练成功 & lane 在当前速率和均衡参数下可传输 & 像素格式、TCON 与面板供电均正确 \\\tablerowrule
扫描输出一帧 & 显示控制器和链路已发送整帧 & 面板响应与背光已经达到目标亮度 \\\tablerowrule
光学传感器读到变化 & 光子已经离开屏幕 & 人眼主观体验不存在拖影或闪烁 \\
\bottomrule
\end{longtable}

\section{十、相机与 ISP：从入射光到可计算图像}

相机是显示通路的逆向互补案例：外界光先经过镜头投到 CMOS 像素阵列，像素把光子积分成电荷，列电路与 ADC 把模拟量变成 Bayer RAW 采样；MIPI CSI-2 再把采样传给 SoC，DMA 写入内存，ISP 完成去马赛克、降噪、颜色和曝光处理。应用收到的图像已经穿过光学、模拟、串行链路、内存和专用计算多个边界。

\begin{pdfdiagram}[x=1cm,y=1cm]
  \node[hw state, minimum width=2.2cm] (scene) at (1.2,4.0) {场景光\\光谱与亮度};
  \node[hw block, minimum width=2.3cm] (lens) at (3.9,4.0) {镜头/光圈\\聚焦与通光量};
  \node[hw memory, minimum width=2.5cm] (sensor) at (6.8,4.0) {CMOS 阵列\\曝光、逐行读出};
  \node[hw compute, minimum width=2.5cm] (adc) at (9.8,4.0) {列放大器/ADC\\Bayer RAW};
  \node[hw block, minimum width=2.3cm] (csi) at (12.7,4.0) {CSI-2 PHY\\分包与串行化};
  \draw[hw bus] (scene) -- (lens);
  \draw[hw bus] (lens) -- (sensor);
  \draw[hw bus] (sensor) -- (adc);
  \draw[hw bus] (adc) -- (csi);

  \node[hw memory, minimum width=2.3cm] (dma) at (12.7,1.1) {DMA/帧缓冲\\行号与时间戳};
  \node[hw compute, minimum width=3.0cm] (isp) at (9.3,1.1) {ISP\\校正、降噪、去马赛克};
  \node[hw control, minimum width=3.0cm] (threea) at (5.4,1.1) {3A 控制\\曝光、白平衡、对焦};
  \node[hw state, minimum width=2.3cm] (image) at (1.6,1.1) {输出图像\\RGB/YUV/统计量};
  \draw[hw bus] (csi) -- (dma);
  \draw[hw bus] (dma) -- (isp);
  \draw[hw bus] (isp) -- (image);
  \draw[hw return] (isp) -- (threea);
  \draw[hw return] (threea) -- (sensor);
\end{pdfdiagram}
\diagramnote{相机数据面与控制面的闭环。上排把光转换成 Bayer RAW 并送入 CSI-2，下排由 DMA 和 ISP 生成图像；ISP 统计量还会驱动自动曝光、自动白平衡和自动对焦，修改下一批帧的传感器与镜头参数。一次控制调整通常要过若干帧才完全可见。}

\topic{CMOS 像素先积分光子，再逐行读出电压}

曝光期间，光电二极管把到达的光子转换为电荷；积分时间越长、模拟增益越高，暗处信号越容易观察，运动模糊、饱和与噪声也会增加。常见彩色传感器在像素上覆盖 Bayer 彩色滤光阵列，每个采样点只测红、绿、蓝中的一种。黑电平、像素响应不均匀、坏点、镜头阴影和读出噪声都在此时进入数据，后续 ISP 只能估计和校正，不能凭空恢复已经饱和或被噪声淹没的信息。

多数移动传感器使用 rolling shutter：各行曝光起止时间不同，快速移动或抖动会把直线拍成斜线。帧时间戳若只记 DMA 完成时刻，就会掩盖第一行与最后一行的实际曝光差；做传感器融合时应保留帧开始、曝光中点或逐行时序。global shutter 同时结束所有像素曝光，运动几何更稳定，但像素电路、噪声、面积和功耗代价不同。

\topic{MIPI CSI-2 负责分包，D-PHY/C-PHY 负责电气传输}

传感器把短包用于帧开始、帧结束和行边界，把长包用于像素负载；包头中的数据类型和虚拟通道可以在同一链路上区分 RAW、嵌入式元数据或多路传感器。CSI-2 定义协议和包，D-PHY 或 C-PHY 定义引脚上的高速传输方式。接收端先检查同步、包头纠错与负载 CRC，再按行写入目标缓冲。

链路错误必须和缓冲错误分开。CRC 或同步错误指向传感器时钟、lane 映射、信号完整性或速率配置；帧开始存在但 DMA 缓冲无完成，通常要查描述符、stride、IOMMU 和中断；连续丢帧还可能是 ISP 或内存带宽不能在下一帧到达前归还缓冲。可靠驱动同时记录帧序号、虚拟通道、错误包、DMA 完成和缓冲代际。

\topic{ISP 是按像素流工作的专用流水线}

典型 ISP 依次执行黑电平校正、坏点修复、镜头阴影校正、降噪、去马赛克、颜色校正、伽马或色调映射、锐化和 RGB/YUV 转换。顺序不能任意交换：例如去马赛克前的噪声仍服从 Bayer 采样结构，过早锐化会把噪声一并放大。高分辨率视频要求流水线每拍接收一个或多个像素，中间结果主要保存在行缓冲而不是整帧往返 DRAM。

自动曝光、自动白平衡和自动对焦构成跨帧反馈。ISP 对当前帧统计亮度直方图、颜色分布和清晰度，控制算法据此更新下一帧或后续若干帧的曝光时间、模拟增益、白平衡增益和镜头位置。参数写入、传感器锁存、实际曝光和 DMA 完成属于不同帧代际；若驱动不标记“参数从第几帧生效”，算法会用旧图像评价新参数，形成振荡。

\begin{longtable}{L{0.17\linewidth}L{0.30\linewidth}L{0.41\linewidth}}
\toprule
异常现象 & 首先检查的边界 & 关键证据 \\
\midrule
整帧过曝且高光剪切 & 曝光与模拟增益 & 曝光寄存器生效帧、RAW 最大码值、直方图 \\\tablerowrule
斜线或果冻形变 & rolling shutter 与运动 & 行时间、曝光区间、陀螺仪时间戳 \\\tablerowrule
彩色噪点或偏色 & RAW/ISP 参数与校准 & Bayer 顺序、黑电平、白平衡增益、颜色矩阵 \\\tablerowrule
随机横线或 CRC 错误 & CSI-2 PHY 与 lane 配置 & 包错误计数、速率、终端匹配、lane 映射 \\\tablerowrule
有帧中断却无应用图像 & DMA、缓冲代际与 ISP 完成 & 描述符、IOMMU fault、帧序号、队列水位 \\
\bottomrule
\end{longtable}

\section{十一、NPU：数据复用、量化与张量完成边界}

NPU 并不是“缩小版 GPU”。它把常见张量算子落实为矩阵乘加阵列、片上 SRAM、搬运引擎和少量向量/标量单元，追求让权重与激活在芯片内尽可能多次复用。峰值 TOPS 只描述乘加阵列在理想数据类型和满利用率下的运算能力；真实性能往往由张量切块、片上容量、外存带宽、量化尺度和不规则算子决定。

\begin{pdfdiagram}[x=1cm,y=1cm]
  \node[hw memory, minimum width=2.4cm] (cmd) at (1.3,3.8) {命令队列\\算子、形状、地址};
  \node[hw control, minimum width=2.4cm] (sched) at (4.2,3.8) {调度/DMA\\切块与预取};
  \node[hw memory, minimum width=2.7cm] (sram) at (7.3,3.8) {片上 SRAM\\激活块与权重块};
  \node[hw compute, minimum width=2.7cm] (array) at (10.6,3.8) {MAC/脉动阵列\\乘法与累加};
  \draw[hw bus] (cmd) -- (sched);
  \draw[hw bus] (sched) -- (sram);
  \draw[hw bus] (sram) -- (array);

  \node[hw compute, minimum width=2.7cm] (post) at (10.6,1.0) {后处理单元\\激活、缩放、量化};
  \node[hw memory, minimum width=2.7cm] (out) at (7.3,1.0) {输出缓冲\\片上或外部内存};
  \node[hw state, minimum width=2.4cm] (done) at (4.2,1.0) {完成项/fence\\状态、代际、错误};
  \node[hw state, minimum width=2.4cm] (host) at (1.3,1.0) {CPU/后续算子\\观察结果};
  \draw[hw bus] (array) -- (post);
  \draw[hw bus] (post) -- (out);
  \draw[hw bus] (out) -- (done);
  \draw[hw bus] (done) -- (host);
\end{pdfdiagram}
\diagramnote{一次 NPU 算子的请求—数据—完成闭环。调度器把大张量切成片上 SRAM 能容纳的块，MAC 阵列反复复用这些数据，后处理单元再完成激活与量化。输出写回、错误状态和 fence 必须按规定顺序发布，峰值 TOPS 不包含这些搬运与完成成本。}

\topic{切块决定同一字节能在片上被使用多少次}

矩阵乘法 $C_{M\times N}=A_{M\times K}B_{K\times N}$ 若每次乘加都从 DRAM 读取两个操作数，外存能耗和带宽会先达到上限。NPU 把 $A$、$B$ 切成 SRAM 可容纳的小块：一个激活块沿多个输出通道复用，一个权重块沿多个输入位置复用，部分和则尽量留在累加器中。脉动阵列让数据在相邻处理单元之间逐拍传递，减少对共享寄存器文件和全局互连的反复访问。

片上 SRAM 容量、阵列尺寸与算子形状若不匹配，利用率会明显下降。矩阵边缘不足一个完整 tile 时，一部分 MAC lane 空闲；深度可分离卷积、归一化、索引和动态形状也未必能填满大阵列。编译器需要选择循环次序、tile 大小、双缓冲和算子融合，使 DMA 在阵列计算当前块时预取下一块，并尽量避免中间张量写回外存。

\topic{量化降低带宽和乘法成本，但引入数值契约}

把 FP32 权重和激活量化为 INT8，可把同样容量容纳的数据量提高四倍，并让阵列集成更多低位宽乘法器。量化值通常满足 $x \approx s(q-z)$，其中 $q$ 是整数码、$s$ 是尺度、$z$ 是零点。矩阵乘加常用更宽的整数累加器保存部分和，再经缩放、舍入、饱和和激活函数转换到输出位宽。

尺度可以按整个张量、按通道或按更细粒度设置。粒度越细，数值误差通常越小，尺度读取和硬件控制越复杂。溢出、舍入模式、有符号性和饱和规则必须由模型转换器、编译器与硬件共同遵守；否则模型可能正常运行却在少数输入上产生系统性偏差。比较 NPU 正确性时，应同时保留浮点参考、各层量化参数和允许误差，而不能只看最终分类是否一致。

\topic{稀疏性只有在编码、调度和阵列都支持时才节省时间}

权重中出现零值并不自动提高性能。硬件必须能用位图、索引或结构化模式跳过零乘法，DMA 也要少搬相应数据；非结构化稀疏的索引和不规则调度可能抵消收益。结构化稀疏把每组元素中非零数量限制为固定模式，较容易映射到规则阵列，但模型训练必须遵守该约束。规格中的“稀疏 TOPS”应与稠密 TOPS 分开理解，并说明所需模式。

\topic{NPU 仍然服从队列、地址、可见性和复位规则}

运行时把算子描述符、张量地址、形状、数据类型与量化参数写入命令队列，执行发布屏障后敲 doorbell。NPU 经 IOMMU 或共享虚拟地址读取缓冲；如果平台并非硬件一致，CPU 与 NPU 之间还要执行 cache 清理和失效。算子完成时，设备必须先让输出写入达到规定的可见点，再更新完成项或 fence；后续 CPU、GPU 或另一个 NPU 队列只能在等待该边界后消费结果。

异常不只有“算错”。描述符非法、张量越界、IOMMU fault、DMA 超时、ECC 错误、看门狗超时和热限频都需要独立状态。复位时应停止新提交，隔离 DMA，标记所有未完成任务失败，重建队列与地址映射，并递增 generation；旧 fence 即使随后到达，也不能被当作新任务完成。由此可见，NPU 与 GPU、NVMe 的计算内容不同，系统边界仍是同一套提交—执行—完成—恢复闭环。

\begin{longtable}{L{0.17\linewidth}L{0.25\linewidth}L{0.29\linewidth}L{0.17\linewidth}}
\toprule
限制来源 & 典型证据 & 优先优化方向 & 不应混淆为 \\
\midrule
阵列利用率低 & 活跃 MAC 比例低、边缘 tile 多 & 调整 tile、批量或算子映射 & 外存一定不足 \\\tablerowrule
外存带宽不足 & DMA 忙、阵列等数、复用率低 & 融合算子、扩大复用、压缩/量化 & 峰值 TOPS 太低 \\\tablerowrule
片上容量不足 & 中间张量频繁溢写 DRAM & 改切块、生命周期与双缓冲 & 计算精度问题 \\\tablerowrule
量化误差过大 & 与浮点参考偏差按层累积 & 校准尺度、提高敏感层位宽 & 驱动完成错误 \\\tablerowrule
队列或地址故障 & fence 不前进、IOMMU/设备 fault & 修复描述符、映射与复位代际 & 模型本身太慢 \\
\bottomrule
\end{longtable}
\section{十二、60Hz 双缓冲案例：渲染完成不等于画面已经可见}

显示通路至少包含三个独立完成域：CPU 把命令交给 GPU，GPU 把新帧写入 back buffer，显示控制器在刷新边界采用该缓冲并逐行扫描。下面用固定 60Hz、双缓冲的案例把这些边界放到同一时间轴。每帧周期约为
\[
  T_{\mathrm{frame}}=\frac{1}{60}\approx16.67\mathrm{ms}.
\]
时刻 0，显示控制器已经开始扫描旧 front buffer。CPU 在 1ms 完成新命令提交，GPU 在 7ms 完成渲染并发布 fence。此时新像素已可供设备读取，但当前扫描周期不能在中途换地址，否则同一画面会混入两代帧。

\begin{pdfdiagram}[x=1cm,y=1cm]
  \node[hw state, minimum width=2.5cm] (submit) at (1.4,4.2) {CPU 提交\\命令与目标缓冲};
  \node[hw compute, minimum width=2.6cm] (render) at (4.5,4.2) {GPU 渲染\\写 back buffer};
  \node[hw state, minimum width=2.6cm] (fence) at (7.7,4.2) {渲染 fence\\像素写入可见};
  \node[hw control, minimum width=2.7cm] (flip) at (10.9,4.2) {VBlank 翻页\\锁存新 plane 地址};
  \node[hw memory, minimum width=2.7cm] (scan) at (10.9,1.2) {显示扫描\\逐行读取新帧};
  \node[hw block, minimum width=2.7cm] (link) at (7.7,1.2) {显示链路/TCON\\传输与行列驱动};
  \node[hw state, minimum width=2.7cm] (pixel) at (4.5,1.2) {面板像素响应\\亮度开始变化};
  \node[hw state, minimum width=2.5cm] (complete) at (1.4,1.2) {整帧可见边界\\最后一行已更新};
  \draw[hw bus] (submit) -- (render);
  \draw[hw bus] (render) -- (fence);
  \draw[hw bus] (fence) -- (flip);
  \draw[hw bus] (flip) -- (scan);
  \draw[hw return] (scan) -- (link);
  \draw[hw return] (link) -- (pixel);
  \draw[hw return] (pixel) -- (complete);
\end{pdfdiagram}
\diagramnote{一帧从命令到可见的四个主要边界。GPU fence 只证明 back buffer 已可读；VBlank 翻页只证明显示控制器采用了新地址；逐行扫描和链路传输之后，面板像素才开始变化；最后一行完成响应，才接近整帧可见。}

\topic{按时间逐项核对这一帧}

在 16.67ms 的 VBlank 边界，显示控制器确认 GPU fence 已完成，锁存 back buffer 地址，并把它变成新的 front buffer。新一轮扫描随后从顶部开始。屏幕顶部像素在 16.67ms 后不久得到新数据，底部像素要接近 33.33ms 才被扫描；LCD/OLED 面板还会增加自身响应时间。因此“这一帧显示了”的精确含义必须说明是首像素变化、整帧扫描完成，还是光学测量达到目标亮度。

\begin{longtable}{L{0.15\linewidth}L{0.24\linewidth}L{0.31\linewidth}L{0.20\linewidth}}
\toprule
时刻 & 事件 & 已经证明的事实 & 尚未证明的事实 \\
\midrule
0ms & 扫描旧 front buffer & 显示控制器正在消费旧帧 & 新帧是否已开始渲染 \\\tablerowrule
1ms & CPU 提交完成 & GPU 可以读取命令与资源描述 & GPU 已经执行或写完像素 \\\tablerowrule
7ms & GPU fence 完成 & back buffer 写入达到设备可见边界 & 当前扫描已采用新缓冲 \\\tablerowrule
16.67ms & VBlank 翻页 & 新 plane 地址成为下一扫描源 & 整个面板都已显示新帧 \\\tablerowrule
16.67--33.33ms & 逐行扫描 & 新帧从顶部向底部送往面板 & 尚未扫描的行仍保持旧代际 \\\tablerowrule
约 33.33ms 以后 & 最后一行及面板响应 & 整帧物理输出接近完成 & 人眼体验仍受响应、背光和处理影响 \\
\bottomrule
\end{longtable}

若 GPU 在 16.90ms 才完成 fence，它已经错过本次 VBlank。显示控制器必须继续扫描旧 front buffer，直到 33.33ms 的下一个边界再翻页。只晚 0.23ms 的计算，可能增加接近一整帧的等待。平均 GPU 时间仍可能低于 16.67ms，但偶发尾延迟会形成明显卡顿，这就是显示系统必须同时观察平均帧率和 frame-time 分布的原因。

\topic{双缓冲、三缓冲与低延迟不是同一个目标}

双缓冲中，未显示的缓冲只有一个。GPU 若完成过早，可能等待显示控制器释放旧 front buffer；若完成过晚，就错过刷新边界。三缓冲允许 GPU 在一个缓冲等待显示时继续渲染另一个缓冲，提高吞吐并减少阻塞，但新帧可能在翻页队列中多等待一代，输入到光子的时延反而增加。低延迟模式通常限制排队深度，让 CPU 不要领先 GPU 和显示过多帧。

可变刷新率改变的是 VBlank 到来的策略：显示器允许在范围内延长帧间隔，GPU 一完成就可以更快触发下一次刷新；一旦一行扫描开始，像素链路仍要按规定速率连续发送。可变刷新不会消除 GPU fence、翻页、扫描和面板响应之间的边界，也不能补救显示带宽不足导致的 underflow。

\begin{longtable}{L{0.18\linewidth}L{0.27\linewidth}L{0.27\linewidth}L{0.18\linewidth}}
\toprule
现象 & 最早检查的边界 & 关键证据 & 恢复或修正 \\
\midrule
命令提交后画面静止 & 命令环与 GPU 执行 & 读指针、fault、fence 值与代际 & 修复队列或重建 GPU 上下文 \\\tablerowrule
fence 完成但仍是旧画面 & 翻页队列与 VBlank & plane 地址、flip 序号、刷新时间戳 & 核对目标缓冲和截止时刻 \\\tablerowrule
画面中部撕裂 & 地址切换时刻 & 翻页是否在扫描中途生效 & 只在允许刷新边界锁存 \\\tablerowrule
随机横线或闪屏 & 扫描带宽与链路 & underflow、链路 CRC/训练、像素时钟 & 增加带宽余量或重新训练 \\\tablerowrule
复位后出现旧帧 & 缓冲与完成代际 & 旧 fence/flip 是否写入新队列 & 递增 generation，丢弃迟到完成 \\
\bottomrule
\end{longtable}

显示恢复也必须遵守所有权顺序。驱动先停止新提交，等待或取消在途 GPU 工作；关闭或冻结 plane 更新；重建显存映射、命令环和显示状态；递增队列 generation；最后恢复渲染与翻页。旧代际 fence 或 flip 完成只能用于诊断，不能释放当前缓冲。这样即使 GPU、显示控制器和面板各自有独立时钟与复位，恢复后也不会把旧帧身份误当作新帧完成。

\section*{章末练习}

\begin{chapterexercise}{综合练习：时延与吞吐取向至占用率演算}
\begin{exercisesubproblem}{（a）时延与吞吐取向}
同一晶体管预算两个方案：A 为 1 个 4 发射乱序核、5 GHz，每拍完成 4 次浮点加；B 为 128 条简单 lane、1.5 GHz，每 lane 每拍 1 次浮点加。对规则数组加法求两方案的峰值吞吐与比值；再说明指针追逐类代码该选哪个、为什么
\exercisecheck{学习目标}{时延取向与吞吐取向的面积分配}
\end{exercisesubproblem}
\begin{exercisesubproblem}{（b）占用率演算}
某 SM 有 65536 个 32 位寄存器，内核每线程用 32 个寄存器，每块 256 线程。求每 SM 最多常驻几块、合多少线程与 warp；若访存时延约 400 拍、每 warp 平均 8 拍发一条指令，估算藏住时延需要的 warp 数并判断该内核够不够；每线程改用 64 个寄存器时又如何
\exercisecheck{学习目标}{在途 warp 数决定时延隐藏}
\end{exercisesubproblem}
\end{chapterexercise}

\begin{chapterexercise}{综合练习：分歧代价至合并访存}
\begin{exercisesubproblem}{（a）分歧代价}
某 warp 执行“为正走 A、否则走 B”的分支，路径 A 约 200 拍、B 约 300 拍，数据随机使约半数 lane 走 A。求该 warp 的耗时，并与全部走 A 的理想情形相比；若 32 条 lane 各走不同分支，吞吐退到多少
\exercisecheck{学习目标}{串行执行两条路径，最坏 1/32}
\end{exercisesubproblem}
\begin{exercisesubproblem}{（b）合并访存}
某 warp 读 32 个连续 float，恰为 128 B 一次事务；改为跨步访问。求跨步 8（\code{a[i*8]}）与跨步 32 时一次 warp 访存各需几次 128 B 事务、带宽效率各为多少；峰值 400 GB/s 下跨步 32 的有效带宽是多少
\exercisecheck{学习目标}{连续地址拼事务，跨步拆事务}
\end{exercisesubproblem}
\end{chapterexercise}

\begin{chapterexercise}{综合练习：带宽预算至内核 trace 排序}
\begin{exercisesubproblem}{（a）带宽预算}
向量加法 $N=2^{20}$（输入 8 MiB、输出 4 MiB），PCIe 3.0 x16 单向 15.75 GB/s，显存带宽 400 GB/s。忽略启动开销，求拷入、内核、拷回各耗时、总时长与内核占比；给出一条降低搬运占比的工程做法
\exercisecheck{学习目标}{算术强度决定值不值得上 GPU}
\end{exercisesubproblem}
\begin{exercisesubproblem}{（b）内核 trace 排序}
把一次内核启动的这些事件排成时间序：写 doorbell、输入 DMA 入显存、fence 更新、启动命令写入命令环、内核在 SM 执行、结果 DMA 回主存、CPU 读结果；并指出哪几处先后是硬性约定
\exercisecheck{学习目标}{提交—执行—完成三段次序}
\end{exercisesubproblem}
\end{chapterexercise}

\begin{chapterexercise}{统一地址}
解释统一虚拟地址仍需页驻留和同步的原因
\exercisecheck{检查要点}{地址命名不等于数据可见}
\end{chapterexercise}

\begin{chapterexercise}{发布顺序}
写出 CPU 准备输入到 GPU fence 的顺序
\exercisecheck{检查要点}{描述符先于 doorbell 可见}
\end{chapterexercise}

\begin{chapterexercise}{一致性}
比较 Cache 一致性与 fence 的职责
\exercisecheck{检查要点}{一个管值，一个管批次完成}
\end{chapterexercise}

\begin{chapterexercise}{页面迁移}
写出一次 GPU 缺页从暂停到重放的步骤
\exercisecheck{检查要点}{更新页表后还要失效设备 TLB}
\end{chapterexercise}

\begin{chapterexercise}{页框回收}
说明 CPU 何时才能释放 GPU 使用过的页面
\exercisecheck{检查要点}{等待完成并撤销设备翻译}
\end{chapterexercise}

\begin{chapterexercise}{显示边界}
区分命令提交、渲染完成和显示完成
\exercisecheck{检查要点}{三个队列/时钟域分别取证}
\end{chapterexercise}

\begin{chapterexercise}{显卡路径}
写出 CPU 提交命令、GPU 读资源、VRAM 写帧、显示扫描的 trace
\exercisecheck{学习目标}{能区分渲染完成、DMA 完成和显示完成}
\end{chapterexercise}

\section*{参考解答与要点}

\begin{chapteranswer}{综合练习：时延与吞吐取向至占用率演算}
\begin{answersubproblem}{（a）时延与吞吐取向}
A：$4\times5\times10^{9}=2\times10^{10}$ 次/秒；B：$128\times1.5\times10^{9}=1.92\times10^{11}$ 次/秒，约为 A 的 9.6 倍。但 B 的单 lane 每拍只 1 次操作且时钟更低，指针追逐这类串行依赖代码在 B 上只有一条 lane 参与执行，时延比 A 差几十倍；数组加法没有依赖，B 的吞吐直接兑现。取向由负载的依赖结构决定
\answercheck{必须掌握的要点}{吞吐方案赢在可并行负载，串行负载仍归时延方案}
\end{answersubproblem}
\begin{answersubproblem}{（b）占用率演算}
每块用 $256\times32=8192$ 个寄存器，$65536/8192=8$，最多 8 块 $=2048$ 线程 $=64$ 条 warp。每 warp 每 8 拍发一条，400 拍时延内需 $400/8=50$ 条可发射 warp；$64>50$，时延可藏住。每线程改用 64 个寄存器后每块占 16384 个，只容 4 块 $=32$ 条 warp，$32<50$，发射槽开始空转，SM 吞吐明显下滑。寄存器用量是占用率的第一约束
\answercheck{必须掌握的要点}{先算每块资源，再算 warp 数，最后对时延需求}
\end{answersubproblem}
\end{chapteranswer}

\begin{chapteranswer}{综合练习：分歧代价至合并访存}
\begin{answersubproblem}{（a）分歧代价}
warp 内两条路径都有 lane 走，硬件串行执行 A 再 B，耗时 $200+300=500$ 拍；全部走 A 只需 200 拍，分歧代价 $500/200=2.5$ 倍。32 条 lane 各走不同分支时 32 条路径全串行，warp 吞吐退到峰值的 $1/32$。分歧的代价由路径数决定，不由走每条路径的 lane 数决定
\answercheck{必须掌握的要点}{按路径串行计费；让同 warp 的 lane 走同一条路}
\end{answersubproblem}
\begin{answersubproblem}{（b）合并访存}
连续：32 个 float 恰为 128 B，1 次事务，效率 1。跨步 8：lane 地址间隔 32 B，warp 覆盖 $32\times32$ B $=1024$ B $=8$ 条线，8 次事务，取 1024 B 只用 128 B，效率 $1/8$。跨步 32：间隔 128 B，每 lane 各碰一条线，32 次事务，效率 $1/32$，有效带宽 $400/32=12.5$ GB/s。访存效率看“一次事务取回的字节里用掉几个”
\answercheck{必须掌握的要点}{事务数等于触及的线数；跨步越大带宽越薄}
\end{answersubproblem}
\end{chapteranswer}

\begin{chapteranswer}{综合练习：带宽预算至内核 trace 排序}
\begin{answersubproblem}{（a）带宽预算}
拷入 $8\times2^{20}/(15.75\times10^{9})\approx533\ \mu$s；内核流量 12 MiB，$12\times2^{20}/(400\times10^{9})\approx31\ \mu$s；拷回 $4\times2^{20}/(15.75\times10^{9})\approx266\ \mu$s。总约 830 $\mu$s，内核占 $31/830\approx3.7\%$，搬运占九成半。做法：数据常驻显存、一次搬运后连发多个内核，或把拷入、内核、拷回分块流水重叠；算术强度这么低的负载，最直接的答案是干脆留在 CPU 上
\answercheck{必须掌握的要点}{PCIe 是瓶颈；摊薄搬运靠批量化与常驻}
\end{answersubproblem}
\begin{answersubproblem}{（b）内核 trace 排序}
次序：输入 DMA 入显存 → 启动命令写入命令环 → 写 doorbell → 内核在 SM 执行 → fence 更新 → 结果 DMA 回主存 → CPU 读结果。三处硬性约定：doorbell 必须在命令可见之后（先可见、再通知）；fence 必须在内核完成之后；CPU 读结果必须在 fence 与回拷完成之后。输入拷贝本身也走“命令加 doorbell”，此处抽象成一步
\answercheck{必须掌握的要点}{与 NVMe 同骨架：提交—执行—完成，次序即可见性}
\end{answersubproblem}
\end{chapteranswer}

\begin{chapteranswer}{统一地址}
两侧可以解释同一地址，但页面可能位于不同介质，Cache 和 TLB 也保存独立状态；仍需迁移、权限和所有权协议。
\answercheck{必须保留的检查点}{指针相同不代表副本最新}
\end{chapteranswer}

\begin{chapteranswer}{发布顺序}
CPU 写数据与描述符，执行所需同步和发布屏障，再写 doorbell；GPU 完成结果写入后更新 fence，CPU 等待并取得可见结果。
\answercheck{必须保留的检查点}{接受命令不等于完成计算}
\end{chapteranswer}

\begin{chapteranswer}{一致性}
一致性维持同一 Cache line 的值关系，fence 确认一批异步工作已经越过完成边界。
\answercheck{必须保留的检查点}{两者不可互相替代}
\end{chapteranswer}

\begin{chapteranswer}{页面迁移}
设备报告 fault 并暂停相关工作，主机分配/复制页面、更新设备页表、失效 TLB，然后重放访问。
\answercheck{必须保留的检查点}{fault 可显著增加尾延迟}
\end{chapteranswer}

\begin{chapteranswer}{页框回收}
阻止新提交，等待相关 fence，撤销并确认设备映射/TLB 失效，确认无在途 DMA 后才能释放。
\answercheck{必须保留的检查点}{不能只改 CPU 页表}
\end{chapteranswer}

\begin{chapteranswer}{显示边界}
doorbell 表示可取命令，GPU fence 表示渲染结果完成，显示完成还要等翻页被采用并扫描输出。
\answercheck{必须保留的检查点}{fence 完成不保证画面已显示}
\end{chapteranswer}

\begin{chapteranswer}{显卡路径}
CPU 写命令环并按 doorbell，GPU 读资源和命令渲染，结果写 VRAM，显示控制器按扫描时序读出；渲染完成、DMA 完成、显示完成是三件不同的事。
\answercheck{必须掌握的要点}{三种完成语义不能混。}
\end{chapteranswer}
